\documentclass[12pt, a4paper]{article}

\usepackage[utf8]{inputenc}
\usepackage[T1]{fontenc}
\usepackage{lmodern}
\usepackage[a4paper, margin=1.6cm]{geometry}
\usepackage[table]{xcolor}
\usepackage{longtable}
\usepackage{array}
\usepackage{booktabs}
\usepackage{enumitem}
\usepackage{titlesec}
\usepackage[most]{tcolorbox}
\usepackage{cite}
\usepackage{hyperref}
\usepackage{fancyhdr}


\usepackage{tabularx}
\usepackage{microtype}
\usepackage{xurl}
\newcommand{\doi}[1]{\href{https://doi.org/#1}{doi:\,\nolinkurl{#1}}}

\definecolor{cNavy}{RGB}{24, 43, 73}
\definecolor{cTerracotta}{RGB}{184, 74, 45}
\definecolor{cTeal}{RGB}{19, 108, 120}
\definecolor{cForest}{RGB}{38, 102, 75}
\definecolor{cPlum}{RGB}{98, 59, 112}
\definecolor{bgSand}{RGB}{251, 248, 243}
\definecolor{bgIce}{RGB}{242, 247, 250}
\definecolor{bgMint}{RGB}{241, 248, 244}
\definecolor{bgRow}{RGB}{247, 249, 251}

\hypersetup{
    colorlinks=true,
    linkcolor=cNavy,
    citecolor=cTerracotta,
    urlcolor=cTeal,
    pdftitle={COMP2300 D1: Interview Analysis, Personas and Scenarios}
}

\titleformat{\section}{\bfseries\color{cNavy}}{\thesection.}{0.5em}{}[\vspace{2pt}{\color{cTerracotta}\titlerule[0.8pt]}]
\titleformat{\subsection}{\bfseries\color{cTerracotta}}{\thesubsection}{0.4em}{}
\titlespacing*{\section}{0pt}{0.8em}{0.32em}
\titlespacing*{\subsection}{0pt}{0.5em}{0.2em}

\pagestyle{fancy}
\fancyhf{}
\fancyhead[L]{\textcolor{gray}{CarbonSaathi D1: Integrated Problem-Space Analysis}}
\fancyhead[R]{\textcolor{gray}{Arnav, Hiranya, Keshav}}
\fancyfoot[C]{\thepage}
\renewcommand{\headrulewidth}{0.4pt}

\newtcolorbox{semanticcard}[1]{
    enhanced,
    colback=bgIce,
    colframe=cTeal,
    boxrule=0pt,
    leftrule=3.5pt,
    arc=2pt,
    left=7pt, right=7pt, top=5pt, bottom=5pt,
    before skip=4.5pt, after skip=4.5pt,
    before upper={\textbf{\textcolor{cNavy}{#1}}\par\vspace{2pt}}
}

\newtcolorbox{latentcard}[1]{
    enhanced,
    colback=bgMint,
    colframe=cForest,
    boxrule=0pt,
    leftrule=3.5pt,
    arc=2pt,
    left=7pt, right=7pt, top=5pt, bottom=5pt,
    before skip=4.5pt, after skip=4.5pt,
    before upper={\textbf{\textcolor{cForest}{#1}}\par\vspace{2pt}}
}

\newtcolorbox{synthesisbox}[1]{
    enhanced,
    colback=bgSand,
    colframe=cTerracotta,
    boxrule=0.8pt,
    leftrule=4pt,
    arc=3pt,
    left=8pt, right=8pt, top=5pt, bottom=5pt,
    before skip=5pt, after skip=5pt,
    before upper={\textbf{\textcolor{cTerracotta}{#1}}\par\vspace{2pt}}
}

\newtcolorbox{scenariocard}[2]{
    enhanced,
    colback=white,
    colframe=#1,
    boxrule=0.7pt,
    leftrule=4pt,
    arc=2pt,
    left=7pt, right=7pt, top=5pt, bottom=5pt,
    before skip=4.5pt, after skip=4.5pt,
    before upper={\textbf{\textcolor{#1}{#2}}\par\vspace{2pt}\hrule\vspace{3pt}}
}

\setlength{\parindent}{0pt}
\setlength{\parskip}{0.26em}
\renewcommand{\arraystretch}{1.15}
\setlist[itemize]{leftmargin=1.3em, itemsep=1pt, topsep=1pt}
\setlist[enumerate]{leftmargin=1.3em, itemsep=1pt, topsep=1pt}




\newtcolorbox{jtbdbox}[1]{
    enhanced,
    colback=bgSand,
    colframe=cTerracotta,
    boxrule=0.8pt,
    leftrule=4pt,
    arc=3pt,
    left=8pt, right=8pt, top=6pt, bottom=6pt,
    before skip=5pt, after skip=5pt,
    before upper={\textbf{\small\textcolor{cTerracotta}{#1}}\par\vspace{2pt}}
}

\newtcolorbox{checkcard}[1]{
    enhanced,
    colback=bgIce,
    colframe=cTeal,
    boxrule=0pt,
    leftrule=3.5pt,
    arc=2pt,
    left=7pt, right=7pt, top=5pt, bottom=5pt,
    before skip=4.5pt, after skip=4.5pt,
    before upper={\textbf{\small\textcolor{cNavy}{#1}}\par\vspace{2pt}}
}
\setlength{\headheight}{15pt}
\begin{document}
\begin{titlepage}
\centering
\vspace*{2.0cm}
{\Large\bfseries\textcolor{cNavy}{COMP2300: Software Design and Development Project}\par}
\vspace{0.8cm}
{\huge\bfseries\textcolor{cNavy}{CarbonSaathi}\par}
\vspace{0.35cm}
{\Large Integrated Problem-Space Analysis\par}
\vspace{0.25cm}
{\large Literature Review, Thematic Analysis, Personas, Scenarios, JTBD and MoSCoW Prioritisation\par}
\vspace{1cm}
{\color{cTerracotta}\rule{0.75\textwidth}{1pt}\par}
\vspace{1cm}
{\large Deliverable 1 (D1) \par}
\vspace{0.4cm}
{\large Topic 1: Carbon Footprint Awareness and Reduction Tool\par}
{\large Urban India Context\par}
\vfill
{\normalsize Arnav Mishra \quad Hiranya Makkar \quad Keshav Singhal\par}
\vspace{0.3cm}
{\normalsize October 2026\par}
\vspace{0.5cm}
{\small Integrated from the team's Literature Review v1.0, Interview Analysis v1.7, and JTBD/Epics v1.9 source documents.\par}
\end{titlepage}

\tableofcontents
\clearpage

\section{Literature Review}
\noindent\textbf{Abstract.} Carbon-footprint tools turn activities into emissions estimates, but a number alone may not help users decide what to change. This review argues that estimates must be legible and recommendations feasible, and identifies gaps for urban Indian students.

\subsection{Introduction}
\subsubsection{Purpose and scope}
Carbon-footprint calculators make the climate impact of transport, energy, food and consumption visible, usually as kilograms of carbon dioxide equivalent (kg\,CO$_2$e), but users also need to understand where an estimate comes from and which actions are realistically open to them. This matters for urban students and young adults, whose choices depend on budgets, shared accommodation and family decisions. The review draws on the brief's starter papers and further peer-reviewed work on calculators, tracking apps and human--computer interaction (HCI). It ends with four design challenges and three research questions (RQ1--RQ3), which the accompanying interview analysis examines through its semantic (S1--S3) and latent (L1--L3) themes.

\subsubsection{Why context matters}
Lee et al.\ link household-consumption surveys to a global supply-chain database and show that footprints vary across places and households in India \cite{lee2021}, so one generic baseline cannot describe every user. Evidence on Indian students is thinner: in a cross-sectional study of 131 Delhi college students, only 38\% said they knew the term ``carbon footprint'' \cite{gandhi2025}, and people estimate carbon quantities poorly relative to familiar metrics such as calories \cite{grinstein2018}. Such surveys measure awareness, not emissions, so they cannot show that an interface changes behaviour. Scope is a second issue: footprints include emissions embodied in goods and services \cite{shi2021}, so estimates should state what they include and adapt categories and units to users.

\subsection{From calculating to enabling change}
Mulrow et al.\ evaluated 31 online calculators and found wide variation in scope, features and regional customisation \cite{mulrow2019}. Among 4,245 users of one widely used calculator, 91\% found it valuable for knowledge and 78\% for motivation, but only 23\% said it gave them the information needed to change their lives \cite{collins2020}.

Tracking-app evidence needs careful reading. Hoffmann et al.\ had 216 German university students log consumption in four domains; emissions computed from their \emph{planned} consumption fell 22.5\% overall, from 12.0\% for mobility to 34.7\% for household activities \cite{hoffmann2024}. These are intended changes over two rounds in a short interval, not observed behaviour, with almost no cross-domain effects; responses also varied with green self-efficacy and self-identity, which supports separate entry points for newcomers and experienced users. Lasarov et al.\ followed 210 participants over three waves and found that goal activation affected emissions differently across domains, with evidence of spillover between them \cite{lasarov2024}, although the two studies share authors and are not independent confirmation. Showing a high-impact activity does not make an alternative affordable, accessible or acceptable, so the design question is how to connect an estimate to an achievable next step.

\begin{table*}[t]
\centering
\footnotesize
\setlength{\tabcolsep}{4pt}
\caption{Design challenges and where the interview analysis examines them.}
\label{tab:challenges}
\begin{tabularx}{\textwidth}{>{\bfseries\raggedright\arraybackslash}p{4.7cm} p{4.5cm} X}
\toprule
\textbf{Design challenge (analysis link)} & \textbf{Evidence base} & \textbf{Implication for the tool}\\
\midrule
Legible, localised estimates (S1, S2; RQ1) & \cite{lee2021,gandhi2025,grinstein2018,mohanty2023,shi2021}: regional variation; unfamiliar terms; mixed evidence on units & Water-bucket and bus-trip comparisons (assumptions stated) and cost in Indian rupees (INR) beside kg\,CO$_2$e; regional grid factors\\
Feasible, non-moralising action (L3) & \cite{collins2020,hoffmann2024,brynjarsdottir2012}: 23\% felt equipped; weak mobility effect & Compare achievable alternatives; mark safety- or work-locked choices; no judgmental scores\\
Low-effort, sustained use (S3; RQ2) & \cite{hoffmann2024,froehlich2010}: effortful tracking; few behaviour studies & Presets, short offline check-ins; simple and detailed views for newcomers and experienced users; language options\\
Shared households (L1, L2; RQ3) & \cite{froehlich2012,tiefenbeck2019}: formative, Western; no Indian evidence & Shared totals before individual shares; no unsupported attribution\\
\bottomrule
\end{tabularx}
\end{table*}

\subsection{Interactive technology and engagement}
HCI research adds a design perspective. A survey of eco-feedback studies found that few HCI papers measured behaviour change \cite{froehlich2010}, so a novel interface is not evidence of effect. In hypothetical ride-hailing surveys, Mohanty et al.\ found that emissions information raised green choices, but raw carbon dioxide (CO$_2$) by weight outperformed relatable equivalencies and people responded to differences between options rather than absolute values \cite{mohanty2023}, so which units work remains unsettled. Brynjarsd\'ottir et al.\ criticise persuasive sustainability work for reducing sustainability to simple metrics \cite{brynjarsdottir2012}, which argues against moralising scores. Household evidence is sparser: formative water eco-feedback work in North American homes probed competition and privacy \cite{froehlich2012}, and real-time shower feedback cut hotel guests' energy use by 11.4\% although they did not pay for it \cite{tiefenbeck2019}, which is relevant to landlord-paid housing. Continued use should not be assumed: tracking is effortful and leaves gaps when users stop entering data \cite{hoffmann2024}, and nothing reviewed shows that streaks or gamification retain users.

\subsection{Synthesis and research gap}
In summary, footprints vary by context \cite{lee2021,gandhi2025}, estimates alone do not equip users to act \cite{mulrow2019,collins2020}, and tracking-app effects differ by domain, with planned change and no Indian sample \cite{hoffmann2024,lasarov2024}. The sources reviewed leave four gaps relevant to the brief: none tests localised units, grid factors or language with Indian students; none tests how a tool should treat choices constrained by safety, shift work, heat or housing; none examines shared or landlord-paid housing in India, where individual bills may not exist; and smart-meter links, cooking fuel and waste are barely covered. These are gaps in the sources reviewed, not in the whole literature. Table~\ref{tab:challenges} turns them into design challenges.

\subsection{Research questions}
\begin{quote}
\small
\textbf{How can an India-specific, personalised carbon-footprint tool help urban Indian students and young adults understand their lifestyle-related emissions and make informed, feasible choices to reduce them?}
\end{quote}
\textbf{RQ1:} Do familiar comparisons (water buckets, bus trips) and cost in INR each improve understanding of an estimate beyond kg\,CO$_2$e alone? \textbf{RQ2:} Which preset inputs reduce logging effort (time, taps) while preserving useful estimates (completeness, comprehension)? \textbf{RQ3:} How can a shared-household view show communal totals without allocating them to individuals?

\clearpage
\section{Thematic Analysis, Personas and Scenarios}
\subsection{Thematic Analysis Methodology}
We analysed the ten interview transcripts (CF1 to CF10) \cite{cf1, cf2, cf3, cf4, cf5, cf6, cf7, cf8, cf9, cf10} using a six-phase reflexive thematic analysis . The interviews cover people aged 19 to 58 living across Pune, Bengaluru, Noida, Chennai and rural Tamil Nadu, Thane and Mumbai, Gurugram, Kolkata, Lucknow, Hyderabad, and Ahmedabad \cite{cf1, cf2, cf3, cf4, cf5, cf6, cf7, cf8, cf9, cf10}. Because the module team provided CF1 to CF10 as illustrative interviews for the coursework brief rather than recorded field interviews \cite{cf1, cf10} , we treat our findings as working design hypotheses that link our literature review to our software requirements. Our literature review focuses on urban Indian students and young adults, whereas the transcripts also include family homes \cite{cf3, cf6, cf8}. We included two family personas (P5 and P6) alongside our student and young-professional personas so that these shared-household habits are properly represented. In our scenario calculations, we use $\text{CO}_2$ with assumed factors, while the literature review uses $\text{kg CO}_2\text{e}$ for footprints in general.

\textbf{Manual Coding Workflow:} We coded all ten transcripts by hand without automated text software so we would not miss everyday domestic details. Following the six phases, we: (1) read through all ten interviews several times to get familiar with the data, (2) wrote first-pass open codes by splitting the initial lead across the three of us (Arnav led CF1 to CF4, Hiranya led CF5 to CF7, and Keshav led CF8 to CF10, before all three of us read and coded the full set), (3) grouped similar open codes into candidate clusters, (4) checked those clusters back against the full set of interviews, (5) built each theme around one main pattern rather than a list of topics, and (6) wrote up the themes below. Instead of treating coding as a mechanical counting exercise, we talked through our notes together to check our own assumptions. For example, discussing CF1 and CF9 helped us see that late-night cab rides were driven by safety and curfews rather than casual convenience \cite{cf1, cf9}.

\subsection{First-Pass Open Coding and Theme Traceability Matrix}

\rowcolors{2}{bgRow}{white}
\begin{longtable}{|>{\raggedright\arraybackslash}p{3.0cm}|>{\raggedright\arraybackslash}p{7.1cm}|>{\raggedright\arraybackslash}p{3.7cm}|>{\raggedright\arraybackslash}p{1.9cm}|}
\hline
\rowcolor{cNavy}
\textcolor{white}{\textbf{Transcript Source}} & \textcolor{white}{\textbf{First-Pass Open Codes Extracted}} & \textcolor{white}{\textbf{Candidate Code Clusters}} & \textcolor{white}{\textbf{Themes}} \\
\hline
\endfirsthead
\hline
\rowcolor{cNavy}
\textcolor{white}{\textbf{Transcript Source}} & \textcolor{white}{\textbf{First-Pass Open Codes Extracted}} & \textcolor{white}{\textbf{Candidate Code Clusters}} & \textcolor{white}{\textbf{Themes}} \\
\hline
\endhead
\hline
\endfoot

\textbf{CF1} \cite{cf1} (Woman, 21, Pune hostel student) & 
15-minute walk to college, weekend bus or auto, cab before 10:30 pm hostel curfew, kg CO2e means nothing, cab rides as a dozen bus trips, veg mess paneer and dal, 3 am exam delivery plastic boxes, kettle and hidden induction plate, shared geyser left on, summer water tanker, bucket bath vs shower, couple of taps and plate photo, Hindi and Marathi, tiny data plan. & 
Everyday physical units, curfew safety need, shared geyser waste, quick low-data input. & 
\textbf{S1, S2, S3, L1, L3} \\
\hline

\textbf{CF2} \cite{cf2} (Man, 29, Bengaluru WFH engineer) & 
Shared company cab Tue and Thu, avoids crowded Metro, split AC 1 pm to 9 pm on WFH days, shocked by electricity bill, bill only shows units and amount, compare to flats like mine in Bengaluru, Swiggy and Zomato foil and plastic, avocado and quinoa vs local food, Amazon box inside a box, fast fashion clothes, order history or card statement, no lecture, smart plug AC cut-off. & 
Unclear bill units, local flat comparison, delivery and shopping waste, non-judgmental smart plug timer. & 
\textbf{S1, S2, S3, L3} \\
\hline

\textbf{CF3} \cite{cf3} (Woman, 46, Noida 3BHK homemaker) & 
Electricity and LPG with backup inverter, 2 ACs running all day in summer, mother-in-law cold vs daughter studying, manages energy alone, ignored by teenagers, pressure cooker dal saves gas, app vegetables in plastic bags going bad in fridge, terrace compost bin vs kabadiwala, twice-weekly water tanker, bucket bath vs 20-minute teen shower, shared family view in Hindi. & 
One-person policing load, generational habits, app grocery spoilage, shared family view. & 
\textbf{S1, S2, S3, L2} \\
\hline

\textbf{CF4} \cite{cf4} (Man, 24, Chennai env.\ sci.\ postgrad) & 
6 km second-hand bicycle, bus or share auto in rain, village two-wheeler, Western emission factors wrong for Indian coal grid, wants to see assumptions, village millets and pulses vs urban processed food, dairy trade-offs without lecturing, village wood and dung cakes vs imported LPG, repair and reuse of clothes and phones, export data for research vs simple view for friends, Tamil, Hindi, and English. & 
Visible grid assumptions, rural and urban cooking fuels, uncredited repair habits, two-level interface. & 
\textbf{S1, S2, S3, L3} \\
\hline

\textbf{CF5} \cite{cf5} (Woman, 34, Thane marketing manager) & 
Thane to Andheri local train, auto or Uber in rain, partner uses scooter, summer night AC, online grocery vans and plastic vs tired Saturday market trip, bubble wrap guilt, balcony compost pot, morning water pressure and long showers, honest easiest habit swap without feeling like a bad person, track together as a couple. & 
Online van vs local market question, delivery packaging guilt, easiest habit swap, couple tracking. & 
\textbf{S1, S2, L2, L3} \\
\hline

\textbf{CF6} \cite{cf6} (Man, 38, Gurugram bank officer) & 
25 km petrol car drive in expressway traffic, wife takes Metro, shared car for school run, 2 split ACs non-stop April to September, LPG rotis and induction, RO water purifier, daughter asks about grocery plastic packets, impact of dropping one car trip a week, show in rupees and kilos, family chart for daughter. & 
Car vs Metro comparison, dual rupee and kilo figures, child curiosity, shared family chart. & 
\textbf{S1, S2, L2} \\
\hline

\textbf{CF7} \cite{cf7} (Woman, 19, Kolkata shared-flat student) & 
Walks to classes, bus, shared auto, or tram with grandmother, Uber too expensive, fans and rice cooker, landlord pays electricity so nobody thinks about it, 2-hour morning water in buckets, corner-shop rolls and biryani, binned leftover food, vegetable stall plastic despite cloth bag, free tool with minimal typing, simple weekly picture, Bengali for roommate. & 
Landlord-paid power, leftover food waste, free low-typing weekly picture, Bengali language access. & 
\textbf{S2, S3, L1} \\
\hline

\textbf{CF8} \cite{cf8} (Man, 58, Lucknow retired official) & 
Scooter and morning walks, misses fortnightly Delhi train, never owned a car, fans and LPG pressure cooker, no AC in June heat, inverter battery trouble, bucket watering for garden, grandson thinks bucket is old-fashioned, gardener mixes separated waste, compost for roses, plain Hindi, cost per bucket and fan hour, grandson can show how it works. & 
Low-waste elder habits, physical unit trust, plain Hindi wording, grandson onboarding help. & 
\textbf{S1, S3, L2} \\
\hline

\textbf{CF9} \cite{cf9} (Woman, 31, Hyderabad night-shift nurse) & 
Night-shift cab because buses have stopped and dark streets are unsafe, day-shift Metro and auto, weekly rota changes, daytime AC from 11 am to sleep, too exhausted to remember lights or cook, hospital canteen packets and delivery plastic, brings steel tiffin when able, work around rota for quiet-day washing, no night notifications while sleeping, one short weekly message. & 
Daytime sleep AC, night transit safety, shift exhaustion, rota-based quiet prompts. & 
\textbf{S2, S3, L3} \\
\hline

\textbf{CF10} \cite{cf10} (Man, 26, Ahmedabad co-living designer) & 
Cycles and takes BRTS, cab in rain feels like a waste, shared lounge AC and common geyser, fan, laptop, and two monitors in room, bill split equally so no pressure to cut back, one-season online t-shirts binned, annual phone upgrade, street food thali and plastic bottles, track shopping and phone upgrades, one monthly challenge such as keeping phone an extra year. & 
Equal bill-split apathy, fast fashion and phone turnover, single monthly challenge. & 
\textbf{S2, S3, L1} \\
\hline
\end{longtable}

\subsection{Semantic Themes (Explicit Patterns Across the Interviews)}

\begin{semanticcard}{Theme S1: Turning Abstract Carbon Figures into Everyday Units and Local Baselines}
\textbf{Main Pattern:} Footprints vary widely across places and households in India \cite{lee2021}. Across age groups, interviewees want to understand how their daily choices matter, yet they lack a familiar vocabulary for abstract carbon units like \textit{``Kg CO2e''} or raw utility bill \textit{``units''} (CF1, CF2, CF6, CF8) \cite{cf1, cf2, cf6, cf8}. In a survey of 131 Delhi college students, only 38\% said they knew the term ``carbon footprint'' \cite{gandhi2025}, and people estimate carbon quantities poorly compared with familiar metrics such as calories \cite{grinstein2018}. In our transcripts, CF4 argues that Western calculators give inaccurate numbers unless they use an Indian coal-heavy grid factor and show their assumptions \cite{cf4}, while CF2 asks to be compared with similar flats in his own city rather than Europe \cite{cf2}. Participants repeatedly ask for impact to be shown in everyday units: equivalent bus rides \cite{cf1}, 15-litre buckets of water \cite{cf1, cf3, cf8}, hourly fan or AC running costs \cite{cf2, cf3, cf8}, and combined Rupee (INR) and kilogram figures \cite{cf6}. In hypothetical ride-hailing surveys, raw $\text{CO}_2$ by weight outperformed relatable equivalencies and people responded to differences between options rather than absolute values \cite{mohanty2023}. We therefore treat interviewees' requests for everyday household units alongside kilogram figures as a stated preference to test in our evaluation.
\begin{itemize}
    \item \textbf{Quotes:} \textit{``If an app said my cab rides for the month were the same as like, a dozen bus trips, I'd probably just take the bus. Give me something I can picture. `Kg CO2e' means nothing to me.''} (CF1 \cite{cf1}). \textit{``A unit of electricity here is much dirtier. We're still heavily on coal... I'd want to see the assumptions.''} (CF4 \cite{cf4}). \textit{``Compare me to flats like mine in Bengaluru.''} (CF2 \cite{cf2}). \textit{``Show me in rupees and kilos, both.''} (CF6 \cite{cf6}). \textit{``Tell me how much one bucket is, and what the fan costs per hour.''} (CF8 \cite{cf8}).
\end{itemize}
\end{semanticcard}

\begin{semanticcard}{Theme S2: Convenience Waste and Fast Replacement vs.\ Uncredited Repair and Reuse}
\textbf{Main Pattern:} Footprints include emissions embodied in goods and services \cite{shi2021}. Most interviewed households eat vegetarian food on most days or have chicken only once or twice a week (CF1, CF3, CF4, CF5, CF6, CF10) \cite{cf1, cf3, cf4, cf5, cf6, cf10}. In the interviews, participants connect their everyday environmental guilt to convenience orders and quick replacements: plastic boxes, foil, and bubble wrap from late-night food and grocery deliveries (CF1, CF2, CF3, CF5, CF6, CF9) \cite{cf1, cf2, cf3, cf5, cf6, cf9}, extra app vegetables going bad in the fridge or binned student takeaway (CF3, CF7) \cite{cf3, cf7}, dairy trade-offs and imported foods versus local seasonal millets (CF1, CF2, CF4) \cite{cf1, cf2, cf4}, and throwing out cheap online t-shirts after one season or upgrading phones every year (CF2, CF10) \cite{cf2, cf10}. Interviewees also point out that a tool should not just count new purchases while ignoring everyday repair and circular habits, such as fixing phones for years and passing down clothes (CF4) \cite{cf4}, selling dry recyclables to the \textit{kabadiwala} (CF3) \cite{cf3}, and composting kitchen scraps (CF3, CF5, CF8) \cite{cf3, cf5, cf8}.
\begin{itemize}
    \item \textbf{Quotes:} \textit{``We order in every night, and then there's just boxes everywhere. I feel bad about the plastic.''} (CF1 \cite{cf1}). \textit{``And the apps, they come in plastic bags, which I hate. And they always send too much, so half of it goes bad in the fridge. That's the waste I really see.''} (CF3 \cite{cf3}). \textit{``I genuinely want to know if the market is better than ordering online, because the van is a factor too.''} (CF5 \cite{cf5}). \textit{``Phone upgrades and clothes are the stuff I'd actually change... keep the phone an extra year.''} (CF10 \cite{cf10}). \textit{``In the village nothing gets thrown away quickly. Clothes get passed down, phones last years, someone fixes everything. If the tool only counts what you buy, it misses all that.''} (CF4 \cite{cf4}).
\end{itemize}
\end{semanticcard}

\begin{semanticcard}{Theme S3: Keeping Logging Effort Low and Offering Simple Multilingual Views}
\textbf{Main Pattern:} Tracking is effortful and leaves gaps when users stop entering data \cite{hoffmann2024}, and a survey of HCI eco-feedback studies found that few papers measured behaviour change \cite{froehlich2010}. In the transcripts, participants treat carbon tracking as a background habit that they will drop if it feels like a long form (CF1, CF2, CF7, CF8, CF9) \cite{cf1, cf2, cf7, cf8, cf9}. Users reject 40-question questionnaires, heavy typing, and high mobile data use (CF1, CF7) \cite{cf1, cf7}, asking instead for a couple of taps, simple weekly or monthly summaries, or a timer on an existing smart plug (CF1, CF2, CF7, CF9, CF10) \cite{cf1, cf2, cf7, cf9, cf10}. Tracking-app responses also varied with green self-efficacy and self-identity, which supports separate entry points for newcomers and experienced users \cite{hoffmann2024}. In our interviews, participants similarly ask for a simple visual view across regional languages (\textbf{Hindi} \cite{cf1, cf3, cf4, cf8}, \textbf{Marathi} \cite{cf1}, \textbf{Tamil} \cite{cf4}, \textbf{Bengali} \cite{cf7}, and \textbf{Gujarati} \cite{cf10}) alongside an optional detailed view where experienced users can check emission assumptions and export CSV data (CF4) \cite{cf4}.
\begin{itemize}
    \item \textbf{Quotes:} \textit{``It can't be a long form. I'm not sitting there doing forty questions. Maybe a couple of taps a day?... And it can't eat data. My plan is tiny.''} (CF1 \cite{cf1}). \textit{``A simple weekly picture would be better than a long report. Oh, and it should be in Bengali too, my roommate doesn't read English well.''} (CF7 \cite{cf7}). \textit{``So maybe a simple view for them and a detailed one for people like me.''} (CF4 \cite{cf4}).
\end{itemize}
\end{semanticcard}

\subsection{Latent Themes (Underlying Household and Structural Dynamics)}

\begin{latentcard}{Theme L1 (Latent): Nobody Wants to Police Shared Bills in Hostels, Flats, and Co-Living}
In university hostels (CF1 \cite{cf1}), rented student flats where the landlord pays the electricity bill (CF7 \cite{cf7}), and co-living spaces that split bills equally (CF10 \cite{cf10}), what one person uses is separated from what they pay. Nobody wants to be the awkward flatmate who polices the shared geyser or lounge AC, and saving power on your own does not lower your rent \cite{cf1, cf7, cf10}. Real-time shower feedback cut hotel guests' energy use by 11.4\% even though they did not pay the bill, which links to landlord-paid housing \cite{tiefenbeck2019}, while formative water eco-feedback work in North American homes probed competition and privacy \cite{froehlich2012}. A tool for shared housing needs to make communal resource use visible without turning flatmates against each other.
\begin{itemize}
    \item \textbf{Quotes:} \textit{``The landlord pays the electricity, so none of us really think about it.''} (CF7 \cite{cf7}). \textit{``The bill gets split between everyone, so there's no real pressure to cut back. Nobody's the one who says turn it off.''} (CF10 \cite{cf10}). \textit{``The geyser is the main thing. Somebody always leaves it on and I'm like, who??''} (CF1 \cite{cf1}).
\end{itemize}
\end{latentcard}

\begin{latentcard}{Theme L2 (Latent): One-Person Policing and Generational Friction in Family Homes}
In joint families and couples (CF3, CF5, CF6, CF8 \cite{cf3, cf5, cf6, cf8}), managing water and power is a family negotiation rather than one person's private choice. Prior household eco-feedback studies are formative and Western, with no Indian evidence on shared households \cite{froehlich2012, tiefenbeck2019}. In CF3, the homemaker carries the tiring job of reminding teenagers about 20-minute showers and lights left on \cite{cf3}. Meanwhile, older relatives practise low-waste habits like bucket baths, pressure-cooker cooking, and fan-only cooling (CF1, CF3, CF8 \cite{cf1, cf3, cf8}), which younger relatives often laugh off as \textit{``old-fashioned''} (CF8 \cite{cf8}). At the same time, children and grandchildren often help out by questioning plastic packets (CF6 \cite{cf6}) or showing elders how phone apps work (CF8 \cite{cf8}).
\begin{itemize}
    \item \textbf{Quotes:} \textit{``I do. Nobody else. Lights off when you leave a room, I'm always saying it... They look at their phones more than they look at me... Everyone should see their own part.''} (CF3 \cite{cf3}). \textit{``I water the garden with a bucket, which my grandson finds very funny, he thinks I'm old-fashioned.''} (CF8 \cite{cf8}). \textit{``We're on the same team, but we don't always agree on what matters.''} (CF5 \cite{cf5}).
\end{itemize}
\end{latentcard}

\begin{latentcard}{Theme L3 (Latent): Safety, Shift Work, and Heat Leaving Users Little Choice}
Among users of one widely used calculator, only 23\% said it gave them the information needed to change their lives \cite{collins2020}, tracking-app feedback showed a weak mobility effect \cite{hoffmann2024}, and persuasive sustainability work has been criticised for reducing sustainability to simple metrics, which argues against moralising scores \cite{brynjarsdottir2012}. In our interviews, several high-emission habits happen because of safety, shift work, or summer heat. Women travelling home before a 10:30 PM hostel curfew or finishing a hospital night shift when buses have stopped take solo cabs to stay safe (CF1, CF9 \cite{cf1, cf9}). Likewise, a nurse sleeping at 11 AM in Hyderabad's heat (CF9 \cite{cf9}) or a student studying until 3 AM in a hostel that bans gas stoves (CF1 \cite{cf1}) cannot simply turn off the AC or cook from scratch. Giving people a bad score for choices they cannot safely avoid only makes them feel lectured and put off (CF2, CF4, CF5, CF9 \cite{cf2, cf4, cf5, cf9}).
\begin{itemize}
    \item \textbf{Quotes:} \textit{``Ola and Uber only when it's late, because the hostel gate shuts at ten thirty and I'm not walking back alone in the dark.''} (CF1 \cite{cf1}). \textit{``Night shift I take a cab, because the buses have stopped... Work around my rota... please, no notifications at night. I'm trying to sleep.''} (CF9 \cite{cf9}). \textit{``Not a lecture please, my mum already does that.''} (CF2 \cite{cf2}). \textit{``Not a score that makes me feel like a bad person. Tell me the easiest thing I can change.''} (CF5 \cite{cf5}).
\end{itemize}
\end{latentcard}

\subsection{Cross-Theme Synthesis and Methodological Reflection}
\begin{synthesisbox}{Synthesis: What is the Design Problem Really About?}
Across S1 to L3, the interviewees show concern about waste and household costs, but they lack a familiar vocabulary for emissions \cite{gandhi2025, grinstein2018}, so the gap is about meaning more than motivation. Existing calculators fall short for these users in three ways. First, under S1 and S2, interviewees ask for an Indian grid assumption \cite{cf4}, visibility into convenience packaging, dairy, clothing, phone upgrades, and repair or composting habits \cite{cf1, cf3, cf4, cf10}, and relatable units (Rupees, 15L buckets, fan-hours, and bus trips) \cite{cf1, cf6, cf8}. Second, under S3 and L3, because night safety, shift work, and heat lock certain habits \cite{cf1, cf2, cf9}, moralising scores \cite{brynjarsdottir2012} need to be replaced by low-effort habit comparisons \cite{cf5, cf6}. Third, under L1 and L2, households and hostels need shared views that balance collective totals with optional individual shares \cite{cf1, cf3, cf7, cf10}.
\vspace{2pt}\hrule\vspace{3pt}
\textbf{Research questions (from the literature review):}
\begin{itemize}
    \item \textbf{RQ1:} Do familiar comparisons (water buckets, bus trips) and cost in INR each improve understanding of an estimate beyond $\text{kg CO}_2\text{e}$ alone?
    \item \textbf{RQ2:} Which preset inputs reduce logging effort (time, taps) while preserving useful estimates (completeness, comprehension)?
    \item \textbf{RQ3:} How can a shared-household view show communal totals without allocating them to individuals?
\end{itemize}
The interview analysis examines these questions through its semantic (S1 to S3) and latent (L1 to L3) themes, matching the four design challenges in Table 1 of the literature review: legible, localised estimates (S1, S2, examined via RQ1), feasible, non-moralising action (L3, which has no separate research question and is treated as a design constraint to test in our final evaluation, as a team decision), low-effort, sustained use (S3, examined via RQ2), and shared households (L1, L2, examined via RQ3). Under RQ3, CF3 asks for family members to \textit{``see their own part''} \cite{cf3}, so our proposed design decision is to show shared totals by default and let each member opt in to showing their own share. That choice goes slightly beyond RQ3 as worded and will be settled in D2.
\vspace{2pt}\hrule\vspace{3pt}
\textbf{Methodological Reflection and MVP Scope:} Because CF1 to CF10 are ten illustrative coursework transcripts \cite{cf1, cf10}, they reflect stated preferences rather than observed long-term behaviour. In our scenarios and prototype calculations, we use an assumed grid factor of $0.716\text{ kg CO}_2/\text{kWh}$ (weighted average including renewables) [VERIFY: confirm the figure and its financial year in CEA v19]. In our MVP, the grid factor will be a visible, editable assumption, which answers CF4's wish to see where numbers come from \cite{cf4}, while whether regional grid factors are offered is left to D2. To keep our Java and JavaFX MVP practical and offline-first, we prioritise local storage, a 3-step setup, 2-tap preset logging, everyday unit conversion, and side-by-side habit comparisons. Features that need external reference tables or file exports are deferred: CSV export (\textbf{Should-Have, Increment 2}) \cite{cf4} and a peer or reference-flat benchmark using a static table (\textbf{Could-Have}) \cite{cf2}.
\end{synthesisbox}

\subsubsection*{Triangulation of Interview Themes with the Literature Review}
\rowcolors{2}{bgRow}{white}
\begin{longtable}{|>{\raggedright\arraybackslash}p{1.2cm}|>{\raggedright\arraybackslash}p{4.5cm}|>{\raggedright\arraybackslash}p{5.5cm}|>{\raggedright\arraybackslash}p{4.3cm}|}
\hline
\rowcolor{cNavy}
\textcolor{white}{\textbf{Theme}} & \textcolor{white}{\textbf{Interview Evidence}} & \textcolor{white}{\textbf{What the LR Says}} & \textcolor{white}{\textbf{Tension or Gap}} \\
\hline
\endfirsthead
\hline
\rowcolor{cNavy}
\textcolor{white}{\textbf{Theme}} & \textcolor{white}{\textbf{Interview Evidence}} & \textcolor{white}{\textbf{What the LR Says}} & \textcolor{white}{\textbf{Tension or Gap}} \\
\hline
\endhead
\hline
\endfoot

\textbf{S1} & Interviewees reject \textit{``Kg CO2e''} and ask for Indian grid assumptions, local flat comparisons, and units like buckets, fan-hours, bus trips, and INR \cite{cf1, cf2, cf4, cf6, cf8}. & Footprints vary across places and households in India \cite{lee2021}. Only 38\% of 131 Delhi students said they knew the term ``carbon footprint'' \textbf{[SOURCE TO VERIFY: the original citation key is missing from the supplied bibliography]}. People estimate carbon quantities poorly compared with familiar metrics \cite{grinstein2018}. Raw $\text{CO}_2$ by weight outperformed relatable equivalencies and people respond to differences between options \cite{mohanty2023}. Calculators vary in regional customisation \cite{mulrow2019}. & \textbf{Tension:} Interviewees ask for tangible equivalents, a stated preference that is untested, and the LR calls the question of units unsettled. \\
\hline

\textbf{S2} & Participants highlight delivery plastic, spoilt app vegetables, dairy vs.\ millets, fast fashion, yearly phone upgrades, repair, composting, and \textit{kabadiwala} sales \cite{cf1, cf2, cf3, cf4, cf5, cf7, cf8, cf10}. & Footprints include emissions embodied in goods and services, so estimates should state what they include \cite{shi2021}. & \textbf{Gap:} The LR says cooking fuel and waste are barely covered in the sources. \\
\hline

\textbf{S3} & Users reject 40-question forms and data-heavy apps, asking for 2-tap presets, simple vs.\ detailed views, and Indian languages \cite{cf1, cf2, cf3, cf4, cf7, cf8, cf9, cf10}. & Tracking is effortful and leaves gaps when users stop entering data \cite{hoffmann2024}. Few HCI eco-feedback papers measured behaviour change \cite{froehlich2010}. Responses varied with green self-efficacy and self-identity, which supports separate entry points for newcomers and experienced users \cite{hoffmann2024}. & \textbf{Gap:} The LR says no source tests language with Indian students, and nothing reviewed shows that streaks or gamification retain users. \\
\hline

\textbf{L1} & Hostel, landlord-paid, and equal-split co-living bills remove financial pressure to turn off shared geysers or lounge ACs \cite{cf1, cf7, cf10}. & Real-time shower feedback cut hotel guests' energy use by 11.4\% although they did not pay, which the LR links to landlord-paid housing \cite{tiefenbeck2019}. Formative water eco-feedback work in North American homes probed competition and privacy \cite{froehlich2012}. & \textbf{Gap:} The LR says no source examines shared or landlord-paid housing in India. \\
\hline

\textbf{L2} & Homemakers police teen showers and lights alone while elders' frugal bucket and fan habits are called old-fashioned \cite{cf1, cf3, cf5, cf6, cf8}. & The same formative household water eco-feedback work \cite{froehlich2012}, which the LR describes as Western with no Indian evidence \cite{froehlich2012, tiefenbeck2019}. & \textbf{Gap:} No Indian evidence on shared households. \\
\hline

\textbf{L3} & Night curfews, shift work, and daytime heat force solo cabs, sleep AC, and exam food delivery, and users reject lecturing scores \cite{cf1, cf2, cf4, cf5, cf9}. & Only 23\% of calculator users said they had the information needed to change \cite{collins2020}. The LR notes a weak mobility effect \cite{hoffmann2024}. Persuasive sustainability work can reduce sustainability to simple metrics, which argues against moralising scores \cite{brynjarsdottir2012}. & \textbf{Gap:} The LR says no source tests how a tool should treat choices constrained by safety, shift work, heat or housing. \\
\hline
\end{longtable}

\subsection{Stakeholder Personas}
We created six personas (P1 to P6) across three stakeholder groups (two personas per group). While the literature review focuses on urban Indian students and young adults (Groups 1 and 2), the transcripts also cover family homes, which is why Group 3 (P5 and P6) is included. Each persona is based on a unique primary transcript (CF1, CF4, CF2, CF9, CF3, CF8), supported by secondary transcripts (CF7, CF10, CF5, CF6) so all ten interviews are covered:

\rowcolors{2}{bgRow}{white}
\begin{longtable}{|>{\raggedright\arraybackslash}p{4.1cm}|>{\raggedright\arraybackslash}p{4.3cm}|>{\raggedright\arraybackslash}p{2.9cm}|>{\raggedright\arraybackslash}p{4.2cm}|}
\hline
\rowcolor{cNavy}
\textcolor{white}{\textbf{Stakeholder Group}} & \textcolor{white}{\textbf{Represented Personas}} & \textcolor{white}{\textbf{Primary (Secondary)}} & \textcolor{white}{\textbf{Associated Scenarios}} \\
\hline
\textbf{Group 1:} Urban Students \& Shared-Housing Youth & \textbf{P1:} Ananya Deshmukh \newline \textbf{P2:} Karthik Sundaram & \textbf{CF1} \cite{cf1} (CF7 \cite{cf7}) \newline \textbf{CF4} \cite{cf4} & \textbf{SC2} (P1: Pune Hostel Floor) \newline \textbf{SC4} (P2: Chennai Co-Living) \\
\hline
\textbf{Group 2:} Time-Poor Working Professionals & \textbf{P3:} Rohan Kulkarni \newline \textbf{P4:} Priya Reddy & \textbf{CF2} \cite{cf2} (CF10 \cite{cf10}) \newline \textbf{CF9} \cite{cf9} (CF5 \cite{cf5}) & \textbf{SC4} (P3: Chennai Co-Living) \newline \textbf{SC3} (P4: Night-Shift Rota) \\
\hline
\textbf{Group 3:} Multi-Generational Household Managers \& Elders & \textbf{P5:} Sunita Sharma \newline \textbf{P6:} Rajeshwar Verma & \textbf{CF3} \cite{cf3} (CF5, CF6 \cite{cf5, cf6}) \newline \textbf{CF8} \cite{cf8} (CF3 \cite{cf3}) & \textbf{SC1} (P5 \& P6: Noida 3BHK Water \& Electricity) \\
\hline
\end{longtable}

\begin{longtable}{|>{\raggedright\arraybackslash}p{2.9cm}|>{\raggedright\arraybackslash}p{13.7cm}|}
\hline
\rowcolor{cTeal}\multicolumn{2}{|p{16.6cm}|}{\textcolor{white}{\textbf{P1: Ananya Deshmukh (Hostel Engineering Student), Group 1, Primary: CF1 \cite{cf1} (Secondary: CF7 \cite{cf7}), Themes: S1, S2, S3, L1, L3}}} \\
\hline
\textbf{Demographics} & Age 21, woman, final-year engineering student living in a Pune hostel with a roommate from Kolkata. Uses a phone with a tiny mobile data plan and reads English, Marathi, and Hindi (her mother prefers Marathi or Hindi, and her roommate prefers Bengali) \cite{cf1, cf7}. \\
\hline
\textbf{Behaviours} & Walks 15 minutes to college on weekdays, takes the PMPML bus or shared auto on weekends, and books Ola or Uber late at night because the hostel gate shuts at 10:30 PM \cite{cf1}. Eats veg mess meals (paneer, dal, rice), makes Maggi or eggs on a kettle and hidden induction plate, orders late-night food delivery in plastic boxes during 3 AM exam study sessions, and bins unfinished takeaway food \cite{cf1, cf7}. Shares a floor geyser and summer tanker water \cite{cf1}. \\
\hline
\textbf{Goals} & (1) Picture how curfew cabs, leftover food, exam plastic boxes, and showers add up in everyday units. (2) Log daily habits in a couple of taps without using mobile data or paying money. (3) Prompt floor-mates to switch off the shared geyser without starting an argument \cite{cf1, cf7}. \\
\hline
\textbf{Attitudes} & Cares about food and plastic waste and feels bad about exam-week delivery boxes, but has to stay practical during late-night exams and curfews \cite{cf1}. \\
\hline
\textbf{Motivations} & Clear visual comparisons (such as cab rides compared with a dozen bus trips, or a bucket bath compared with a shower) and a simple weekly picture \cite{cf1, cf7}. \\
\hline
\textbf{Frustrations} & ``Kg CO2e'' means nothing to her, fixed hostel and landlord bills hide electricity costs so people leave the geyser on, and she refuses to fill in 40-question forms \cite{cf1, cf7}. \\
\hline
\end{longtable}

\begin{longtable}{|>{\raggedright\arraybackslash}p{2.9cm}|>{\raggedright\arraybackslash}p{13.7cm}|}
\hline
\rowcolor{cTeal}\multicolumn{2}{|p{16.6cm}|}{\textcolor{white}{\textbf{P2: Karthik Sundaram (Env.\ Sci.\ Postgrad \& Repair Advocate), Group 1, Primary: CF4 \cite{cf4} (Secondary: None), Themes: S1, S2, S3, L3}}} \\
\hline
\textbf{Demographics} & Age 24, man, environmental science postgraduate living in shared accommodation in Chennai after growing up in a village in Tamil Nadu. Uses a laptop for research data and speaks Tamil, Hindi, and English \cite{cf4}. \\
\hline
\textbf{Behaviours} & Cycles 6 km to university on a second-hand bicycle, takes the bus or a share auto when it rains, and uses his father's two-wheeler in his home village \cite{cf4}. Eats vegetarian meals with local seasonal millets and pulses, keeps phones repaired for years, and wears passed-down clothes \cite{cf4}. \\
\hline
\textbf{Goals} & (1) Check the calculator's assumptions to verify it uses an Indian coal-grid factor and accounts for village vs.\ city cooking fuels (wood and dung cakes vs.\ LPG). (2) Use CSV export (\textbf{Should-Have, Increment 2}) for his research while sharing a Simple View in Tamil or Hindi with friends. (3) Get credit for repairing items and cycling \cite{cf4}. \\
\hline
\textbf{Attitudes} & Analytical, sceptical of Western calculators, and against lecturing users \cite{cf4}. \\
\hline
\textbf{Motivations} & Transparent emission assumptions, fair explanations of dairy trade-offs, and recognition for repair and reuse habits \cite{cf4}. \\
\hline
\textbf{Frustrations} & Believes Western tools underestimate India's coal-heavy grid and only count new purchases while ignoring repair culture \cite{cf4}. \\
\hline
\end{longtable}

\begin{longtable}{|>{\raggedright\arraybackslash}p{2.9cm}|>{\raggedright\arraybackslash}p{13.7cm}|}
\hline
\rowcolor{cNavy}\multicolumn{2}{|p{16.6cm}|}{\textcolor{white}{\textbf{P3: Rohan Kulkarni (Hybrid Tech Engineer \& Fast-Upgrade Shopper), Group 2, Primary: CF2 \cite{cf2} (Secondary: CF10 \cite{cf10}), Themes: S1, S2, S3, L1, L3}}} \\
\hline
\textbf{Demographics} & Age 29, man, hybrid software engineer sharing a 2BHK co-living flat in Chennai with Karthik (P2). Speaks English, Gujarati, and Hindi. \textit{Composite Note:} Combines CF2's hybrid WFH AC and delivery habits with CF10's shared co-living bill split and yearly phone upgrades, relocated from Bengaluru and Ahmedabad to Chennai so P2 and P3 share one flat in SC4. The CF2 request to compare against similar flats in Bengaluru is kept as a need for a city-matched reference rather than a Bengaluru-specific one \cite{cf2, cf10}. \\
\hline
\textbf{Behaviours} & Takes a shared company cab to the office on Tuesdays and Thursdays, and works from home the rest of the week running his bedroom split AC from 1 PM to 9 PM alongside two monitors, a laptop, a small fridge, and an air purifier while splitting the flat's electricity bill equally \cite{cf2, cf10}. Orders Swiggy or Zomato two or three times a week, buys avocado and quinoa, upgrades his phone every year, and bins cheap online t-shirts after one season \cite{cf2, cf10}. \\
\hline
\textbf{Goals} & (1) Opt in to seeing his room's share of the split electricity bill and compare against a peer or reference-flat benchmark using a static table (\textbf{Could-Have}). (2) Track fast fashion, yearly phone upgrades, and delivery packaging without daily typing. (3) Set a smart-plug AC timer and take one monthly challenge (keep his phone an extra year) \cite{cf2, cf10}. \\
\hline
\textbf{Attitudes} & Curious after a high monthly electricity bill, but refuses daily logging or lectures \cite{cf2, cf10}. \\
\hline
\textbf{Motivations} & One honest monthly figure, clear room-level AC estimates, and a single doable monthly challenge \cite{cf2, cf10}. \\
\hline
\textbf{Frustrations} & Utility bills only show raw ``units'' and an amount, equal bill-splitting removes pressure to cut back on shared lounge AC, and Amazon and Swiggy orders create heavy packaging waste \cite{cf2, cf10}. \\
\hline
\end{longtable}

\begin{longtable}{|>{\raggedright\arraybackslash}p{2.9cm}|>{\raggedright\arraybackslash}p{13.7cm}|}
\hline
\rowcolor{cNavy}\multicolumn{2}{|p{16.6cm}|}{\textcolor{white}{\textbf{P4: Priya Reddy (Rotating Night-Shift Staff Nurse), Group 2, Primary: CF9 \cite{cf9} (Secondary: CF5 \cite{cf5}), Themes: S2, S3, L3}}} \\
\hline
\textbf{Demographics} & Age 31, woman, hospital staff nurse on rotating night shifts living alone in a flat in Hyderabad. Uses her phone for short weekly updates and reads English and Hindi \cite{cf9}. \\
\hline
\textbf{Behaviours} & Takes the Metro and an auto on day shifts, but has to take a cab on night shifts because the buses have stopped and walking alone in the dark is unsafe \cite{cf9}. Sleeps during the day with her bedroom AC on from 11 AM \cite{cf9}. Exhaustion makes her forget to turn lights off, charge her phone and laptop overnight, and rely on hospital canteen packets and food delivery, though she brings a steel tiffin on days she can manage it \cite{cf9}. \\
\hline
\textbf{Goals} & (1) Have the tool work around her weekly shift rota without penalising night safety cabs or daytime sleep AC. (2) Get one short weekly message (with sleep-time notifications muted) pointing out a quiet day off for washing, tiffin prep, or local market shopping \cite{cf5, cf9}. \\
\hline
\textbf{Attitudes} & Bothered by hospital canteen and delivery plastic waste, but worn out by rotating night shifts \cite{cf9}. \\
\hline
\textbf{Motivations} & Rota-aware reminders and practical, guilt-free suggestions for her quiet days \cite{cf5, cf9}. \\
\hline
\textbf{Frustrations} & Notifications that wake her up while sleeping, scores that make her feel like a bad person, and a shift rota that changes every week \cite{cf5, cf9}. \\
\hline
\end{longtable}

\begin{longtable}{|>{\raggedright\arraybackslash}p{2.9cm}|>{\raggedright\arraybackslash}p{13.7cm}|}
\hline
\rowcolor{cTerracotta}\multicolumn{2}{|p{16.6cm}|}{\textcolor{white}{\textbf{P5: Sunita Sharma (Multi-Generational Household Manager), Group 3, Primary: CF3 \cite{cf3} (Secondary: CF5, CF6 \cite{cf5, cf6}), Themes: S1, S2, S3, L2}}} \\
\hline
\textbf{Demographics} & Age 46, woman, homemaker and part-time accounts assistant in a Noida 3BHK flat (living with her bank-officer husband who drives 25 km to Gurugram, two teenagers, her mother-in-law, and her father-in-law Rajeshwar P6). Uses her phone, the inverter display, and one geyser smart plug, and reads English and Hindi \cite{cf3, cf6, cf8}. \\
\hline
\textbf{Behaviours} & Manages all home energy, water, and waste on her own: cooks dal in an LPG pressure cooker to save gas, does washing in the afternoon, puts the morning geyser on a timer, composts wet waste on the terrace, and sells dry waste to the \textit{kabadiwala} \cite{cf3}. Orders vegetables on apps (which arrive in plastic bags and often go bad in the fridge) and relies on twice-weekly society water tankers \cite{cf3, cf6}. \\
\hline
\textbf{Goals} & (1) Show shared household totals by default while letting each family member decide for themselves whether to opt in to showing their own share, so she is not the only one checking. (2) Show teen showers in litres or buckets and commute or grocery swaps in Rupees and kilos. (3) Track how much waste the house composts and sells to the \textit{kabadiwala} versus throws away \cite{cf3, cf5, cf6}. \\
\hline
\textbf{Attitudes} & Thrifty and practical, but tired of being the only person reminding the house to save power and water \cite{cf3}. \\
\hline
\textbf{Motivations} & Phone charts that her teenagers and husband will actually check, while showing respect for her in-laws' bucket baths \cite{cf3, cf6}. \\
\hline
\textbf{Frustrations} & Teenagers ignore her reminders about 20-minute showers, unattended TVs, and overnight laptop charging, family members disagree over AC use, and app vegetables rot in the fridge \cite{cf3}. \\
\hline
\end{longtable}

\begin{longtable}{|>{\raggedright\arraybackslash}p{2.9cm}|>{\raggedright\arraybackslash}p{13.7cm}|}
\hline
\rowcolor{cTerracotta}\multicolumn{2}{|p{16.6cm}|}{\textcolor{white}{\textbf{P6: Rajeshwar Prasad Verma (Frugal Elder \& Conservationist), Group 3, Primary: CF8 \cite{cf8} (Secondary: CF3 \cite{cf3}), Themes: S1, S3, L2}}} \\
\hline
\textbf{Demographics} & Age 58, man, retired railway official and Sunita's (P5) father-in-law, living with the family in the Noida 3BHK flat (so Sunita's two teenagers are his grandchildren). Dislikes complicated menus, needs Plain Hindi, and asks his grandchildren to show him how apps work. \textit{Composite Note:} His Lucknow details from CF8 (where CF8 mentions one grandson) were relocated to Noida so he can appear as Sunita's father-in-law in SC1 \cite{cf3, cf8}. \\
\hline
\textbf{Behaviours} & Walks in the mornings and rides a scooter to the market and temple, never owned a car, and misses fortnightly train trips to Delhi \cite{cf8}. Uses ceiling fans, an LPG pressure cooker, and a backup inverter, choosing no AC even in June heat \cite{cf3, cf8}. Uses a bucket for bathing and watering his roses, and composts kitchen scraps \cite{cf3, cf8}. \\
\hline
\textbf{Goals} & (1) Check household water and electricity in plain Hindi using familiar units (cost per hour of fan use and buckets of water). (2) Use the app with his grandchildren so they see the value of his low-waste habits \cite{cf3, cf8}. \\
\hline
\textbf{Attitudes} & Frugal, proud of his compost-fed roses, and distrustful of numbers he cannot understand \cite{cf8}. \\
\hline
\textbf{Motivations} & Plain Hindi wording, clear rupee and bucket figures, and spending time with his grandchildren \cite{cf8}. \\
\hline
\textbf{Frustrations} & Grandchildren find his bucket watering funny and call him ``old-fashioned'', colony staff mix up separated waste, and English jargon puts him off \cite{cf8}. \\
\hline
\end{longtable}

\subsection{User Scenarios}

\begin{synthesisbox}{Assumed values used in the scenarios}
In line with Table 1 of the literature review, which asks for water-bucket and bus-trip comparisons to state their assumptions, every number below that is not in a transcript is marked as an illustrative assumption [TEAM: replace with a cited source only if you find one you trust]:
\begin{itemize}
    \item \textbf{Grid factor:} $0.716\text{ kg CO}_2/\text{kWh}$ (assumed, illustrative, not sourced) [VERIFY: confirm the figure and its financial year in CEA v19]
    \item \textbf{Bus factor:} $0.015\text{ kg CO}_2$ per passenger-km (assumed, illustrative, not sourced) [VERIFY: confirm the factor in the source]
    \item \textbf{Student bus round trip:} $15\text{ km}$ assumed student commute (assumed, illustrative, not sourced)
    \item \textbf{Shower flow rate:} $7.5\text{ L/min}$ (assumed, illustrative, not sourced)
    \item \textbf{Electricity tariff:} $\text{Rs.\ } 8\text{ per kWh}$ (assumed, illustrative, not sourced)
    \item \textbf{Ceiling fan power:} $75\text{ W}$ ($0.075\text{ kW}$) (assumed, illustrative, not sourced)
    \item \textbf{Split AC power:} $1.5\text{ kW}$ (assumed, illustrative, not sourced)
    \item \textbf{Bathroom geyser power:} $2\text{ kW}$ (assumed, illustrative, not sourced)
    \item \textbf{Monthly WFH days:} $14\text{ WFH days a month}$ (assumed, illustrative, not sourced)
    \item \textbf{Bucket capacity:} $15\text{ L}$ bucket (assumed, illustrative, not sourced)
\end{itemize}
\end{synthesisbox}

\begin{scenariocard}{cTerracotta}{Scenario 1 (SC1): Resolving Summer Tanker-Water and Fan/Geyser Friction in a Noida 3BHK Home}
\begin{itemize}
    \item \textbf{Stakeholders \& Personas:} \textbf{Group 3} (\textbf{P5 Sunita}, her father-in-law \textbf{P6 Rajeshwar} living with the family in Noida) and Household Teenagers (Rajeshwar's grandchildren), covering Themes S1, S2, S3, and L2.
    \item \textbf{Specific Situation:} On a hot June morning in the Noida 3BHK flat, society water tankers arrive only twice a week \cite{cf3}. Sunita (\textbf{P5}) is tired of reminding her two teenagers about their 20-minute morning showers and unattended fans and TVs \cite{cf3}, while her father-in-law Rajeshwar (\textbf{P6}) wants to see household water and fan costs in plain Hindi \cite{cf8}.
    \item \textbf{Direct System Use:} Sunita opens the \textbf{Household View}, which shows the shared household total by default, logs the morning geyser smart-plug timer and monthly electricity units, and switches the display to \textbf{Bilingual Mode (English and Plain Hindi)} \cite{cf3}. Each family member decides for themselves on their own profile whether to opt in to showing their individual share, and Sunita does not switch attribution on for other people. If a teenager chooses not to opt in, the screen shows the shared household total only, with no pressure message.
    \item \textbf{Indirect Stakeholder Use:} Looking at the screen with his grandchildren, Rajeshwar (\textbf{P6}, assisted user) reads the plain Hindi figures (\textbf{1 ceiling fan at $75\text{ W} = \text{Rs.\ } 0.60/\text{hour}$} at Rs.\ 8/kWh, and $0.075\text{ kWh} \times 0.716 = 0.054\text{ kg CO}_2/\text{hour}$) and sees his bucket watering and terrace composting credited \cite{cf3, cf8}. When Sunita's teenagers check the water comparison on their phones, they see that one 20-minute shower ($150\text{ L}$ at an assumed $7.5\text{ L/min}$) uses as much scarce tanker water as \textbf{ten 15-litre bucket baths} \cite{cf3}.
    \item \textbf{Goal Achieved:} Sunita no longer has to police the bathrooms on her own: the teenagers shorten their morning showers to save tanker water, and Rajeshwar's bucket and fan habits are respected in plain Hindi \cite{cf3, cf8}.
\end{itemize}
\end{scenariocard}

\begin{scenariocard}{cTeal}{Scenario 2 (SC2): Late-Night Exam Delivery Logging and Shared Geyser Nudges in a Pune Hostel}
\begin{itemize}
    \item \textbf{Stakeholders \& Personas:} \textbf{Group 1} (\textbf{P1 Ananya} and her Pune Hostel Floor-Mates), covering Themes S1, S2, S3, L1, and L3.
    \item \textbf{Specific Situation:} Studying until 3 AM during final-year engineering exams in her Pune hostel, Ananya (\textbf{P1}) orders late-night food delivery in plastic boxes, makes Maggi on a hidden induction plate, and notices someone on her floor left the shared bathroom geyser switched on for 2 wasted hours \cite{cf1, cf7}.
    \item \textbf{Direct System Use:} Using offline \textbf{2-Tap Preset Logging} (zero mobile data), Ananya logs her delivery meal, plastic packaging, and induction use in 5 seconds, while her late-night pre-curfew Ola cab ride is marked as a \textit{Safety-Locked Trip} so she is not penalised \cite{cf1}. She then taps the shared-appliance calculator to check the impact of leaving an assumed $2\text{ kW}$ floor geyser on for 2 wasted hours ($4\text{ kWh} \times 0.716\text{ kg CO}_2/\text{kWh} = 2.86\text{ kg CO}_2$) \cite{cf1}.
    \item \textbf{Indirect Stakeholder Use:} Using the assumed bus factor of $0.015\text{ kg CO}_2$ per passenger-km across an assumed $15\text{ km}$ student bus round trip ($0.225\text{ kg CO}_2/\text{trip}$), the tool translates the $2.86\text{ kg CO}_2$ geyser waste into \textbf{13 city bus trips}. Ananya shares the Marathi, Hindi, and Bengali card in her hostel floor chat so floor-mates (indirect users) switch the geyser off \cite{cf1, cf7}.
    \item \textbf{Goal Achieved:} Ananya logs her exam-week packaging waste in seconds without data drain or curfew guilt, and her floor-mates switch off the shared geyser without corridor arguments \cite{cf1}.
\end{itemize}
\end{scenariocard}

\begin{scenariocard}{cNavy}{Scenario 3 (SC3): Night-Shift Rota Quiet Hours and Rest-Day Planning in a Hyderabad Flat}
\begin{itemize}
    \item \textbf{Stakeholders \& Personas:} \textbf{Group 2} (\textbf{P4 Priya}), covering Themes S2, S3, and L3.
    \item \textbf{Specific Situation:} Priya (\textbf{P4}) finishes hospital night shifts in Hyderabad. Because public buses have stopped running at night, she takes solo cabs home to stay safe, orders takeaway in plastic packets because she is too tired to cook, and runs her bedroom AC from 11 AM so she can sleep in the daytime heat \cite{cf9}.
    \item \textbf{Direct System Use:} At the start of the week, Priya enters her rotating shift rota in two taps \cite{cf9}. The tool mutes all alerts during her daytime sleep hours and marks her night-shift cabs and daytime sleep AC as \textit{Health and Safety Necessities}, avoiding any negative score penalty \cite{cf5, cf9}.
    \item \textbf{Indirect / Scheduled Follow-Through:} On her next quiet rest day on the rota, the tool sends \textbf{one short weekly message} suggesting two realistic tasks for her afternoon: running her washing-machine load and packing her reusable steel tiffin for the upcoming shift to avoid single-use hospital canteen packets \cite{cf9}.
    \item \textbf{Goal Achieved:} Priya cuts down on canteen plastic and organises laundry on her quiet days without being woken up during the day or made to feel guilty about her night shifts \cite{cf9}.
\end{itemize}
\end{scenariocard}

\begin{scenariocard}{cPlum}{Scenario 4 (SC4): Shared Co-Living Bill Audit, WFH AC, and Phone Longevity in a Chennai 2BHK}
\begin{itemize}
    \item \textbf{Stakeholders \& Personas:} \textbf{Group 1} (\textbf{P2 Karthik}) and \textbf{Group 2} (\textbf{P3 Rohan}), covering Themes S1, S2, S3, L1, and L3.
    \item \textbf{Specific Situation:} Sharing a Chennai 2BHK co-living flat where electricity bills are split equally, environmental science postgrad Karthik (\textbf{P2}) and hybrid software engineer Rohan (\textbf{P3}) receive a high summer electricity bill after Rohan runs his bedroom's assumed $1.5\text{ kW}$ split AC from 1 PM to 9 PM across three weekly WFH days alongside dual monitors \cite{cf2, cf4, cf10}.
    \item \textbf{Direct System Use:} The flat dashboard shows the shared total by default, and both flatmates choose to opt in to room-level attribution. Karthik opens the \textbf{Detailed Audit View} to inspect the visible, editable grid factor ($0.716\text{ kg CO}_2/\text{kWh}$) \cite{cf4}, logs his bicycle commute and phone repair for circularity credit, and uses CSV export (\textbf{Should-Have, Increment 2}) for his research \cite{cf4}. Rohan checks the \textbf{Simple View}, compares his usage against a peer or reference-flat benchmark using a static table (\textbf{Could-Have}), and tests cutting his room AC off 1 hour earlier across his assumed 14 monthly WFH days ($1.5\text{ kW} \times 1\text{ hr} \times 14\text{ days} = 21\text{ kWh} \times 0.716\text{ kg CO}_2/\text{kWh} = 15.04\text{ kg CO}_2$ and $\text{Rs.\ } 168/\text{month}$ saved at an assumed Rs.\ 8/kWh) \cite{cf2}.
    \item \textbf{Indirect / Monthly Challenge Follow-Through:} Seeing his online shopping summary, Rohan also picks one \textbf{Monthly Challenge}: \textit{``Keep your current smartphone 1 extra year''} \cite{cf10}, while setting his smart-plug timer to turn his room AC off 1 hour earlier automatically \cite{cf2}.
    \item \textbf{Goal Achieved:} Both flatmates see a transparent breakdown of their shared bill without conflict, Rohan automates his WFH AC cut-off and keeps his phone an extra year, and Karthik verifies the calculation assumptions \cite{cf2, cf4, cf10}.
\end{itemize}
\end{scenariocard}




\subsection*{Project management techniques used}
We used a simple Scrum workflow fitted for a three-person research phase. We held two group meetings a week (one in-person whiteboard session and one online catch-up) to review drafts and assign tasks with internal deadlines set 48 hours early. Every section was peer-reviewed by another team member before merging so our research questions and citations stayed consistent.




\subsection*{Reflection on the process (what went well, what to improve)}
Splitting the ten transcripts into three lead batches and second-coding each other's notes worked well and helped us spot our latent themes around shared bills (L1) and night time safety (L3). However, drafting the literature review and interview analysis in separate files initially caused slight differences in our research question wording and persona locations. For Deliverable 2, we will work from a single shared document of personas and epics from day one and track all sprint tasks on GitLab Issues.

\clearpage
\section{Jobs to Be Done, Epics and MoSCoW Prioritisation}
\subsection{Scope and Evidence Base}
This section turns our problem-space findings into a single stakeholder goal (the Job to Be Done, JTBD) and a MoSCoW-prioritised set of epics. It builds directly on our companion Deliverable 1 sections: the literature review \cite{lee2021,gandhi2025,grinstein2018,shi2021,mulrow2019,collins2020,hoffmann2024,lasarov2024,froehlich2010,mohanty2023,brynjarsdottir2012,froehlich2012,tiefenbeck2019} and the interview analysis (version 1.7) of the ten coursework transcripts (\texttt{CF1} to \texttt{CF10}) \cite{cf1,cf2,cf3,cf4,cf5,cf6,cf7,cf8,cf9,cf10}, whose semantic themes (\texttt{S1}--\texttt{S3}), latent themes (\texttt{L1}--\texttt{L3}), personas (\texttt{P1}--\texttt{P6}), scenarios (\texttt{SC1}--\texttt{SC4}), and research questions (\texttt{RQ1}--\texttt{RQ3}) are referenced by ID throughout. An epic is a high-level statement of stakeholder value rather than a technical specification. Epics will be split into user stories with testable acceptance criteria in later deliverables. Because \texttt{CF1} to \texttt{CF10} are illustrative coursework materials rather than a representative population sample \cite{cf1,cf10}, the priorities below are reasoned design judgements grounded in our evidence, not measured demand.

% =========================================================================
% SECTION 2: JOB TO BE DONE
% =========================================================================
\subsection{Job to Be Done (JTBD)}

\begin{jtbdbox}{JTBD Statement (Exact Word Count: 62 Words --- Strictly Under the 100-Word Limit)}
\textit{When I want to understand how my everyday routines and household habits affect the environment, I want to see which activities matter most in terms I can picture and which realistic alternatives are open to me, so that I can make changes that fit my budget, safety, rota, and shared home without being misled, lectured, or blamed for constraints I cannot control.}
\end{jtbdbox}

The statement describes the outcome a stakeholder is hiring a tool to achieve and names no interface, software feature, metric, or technology. It stops at informed, feasible choices rather than guaranteed emission cuts because the reviewed tracking-app studies report intended rather than observed change \cite{hoffmann2024,lasarov2024} and few eco-feedback papers measure behavioural outcomes \cite{froehlich2010}.

\subsubsection{The Job in Each Stakeholder's Terms}
The single JTBD is deliberately broad enough to cover distinct stakeholder circumstances without turning into a feature list. The table below makes that coverage explicit; persona labels follow Interview Analysis v1.7.
{\small
\rowcolors{2}{bgRow}{white}
\begin{longtable}{|>{\raggedright\arraybackslash}p{2.1cm}|>{\raggedright\arraybackslash}p{3.0cm}|>{\raggedright\arraybackslash}p{12.3cm}|}
\hline
\rowcolor{cNavy}
\textcolor{white}{\textbf{Persona}} & \textcolor{white}{\textbf{Core circumstance}} & \textcolor{white}{\textbf{What success means for this stakeholder}} \\
\hline
\endfirsthead
\hline
\rowcolor{cNavy}
\textcolor{white}{\textbf{Persona}} & \textcolor{white}{\textbf{Core circumstance}} & \textcolor{white}{\textbf{What success means for this stakeholder}} \\
\hline
\endhead
\textbf{P1 Ananya} & Student in shared/hostel housing & Understand a few high-impact daily choices without assuming control over fixed bills, curfews, or late-night safety. \\
\hline
\textbf{P2 Karthik} & Evidence-conscious shared-household user & Inspect assumptions and distinguish sourced factors from placeholders rather than trusting a black-box score. \\
\hline
\textbf{P3 Rohan} & Working professional in a shared flat & Compare practical travel and home-energy alternatives in terms of likely impact and cost. \\
\hline
\textbf{P4 Priya} & Rotating night-shift worker & Receive advice that respects safety, heat, work rota, and daytime sleep rather than treating every lower-carbon option as feasible. \\
\hline
\textbf{P5 Sunita} & Household manager & Recognise household routines and low-waste practices without making one person responsible for policing everyone else. \\
\hline
\textbf{P6 Rajeshwar} & Older household member & Understand the result through familiar language and everyday examples without being judged for established routines. \\
\hline
\end{longtable}
}

\subsubsection{Forces Shaping the Job}
{\footnotesize
\rowcolors{2}{bgRow}{white}
\begin{longtable}{|>{\raggedright\arraybackslash}p{2.5cm}|>{\raggedright\arraybackslash}p{9.2cm}|>{\raggedright\arraybackslash}p{4.3cm}|}
\hline
\rowcolor{cNavy}
\textcolor{white}{\textbf{Force}} & \textcolor{white}{\textbf{What the Evidence Shows}} & \textcolor{white}{\textbf{Source (TA \& LR)}} \\
\hline
\textbf{Push}\newline (why current tools fail) & Footprints vary across places and households in India \cite{lee2021}, and calculators vary in scope, features, and regional customisation \cite{mulrow2019}. Only 38\% of 131 Delhi students knew the term ``carbon footprint'' \cite{gandhi2025}, and people estimate carbon poorly relative to familiar metrics such as calories \cite{grinstein2018}. Raw bill ``units'' and carbon weights mean nothing to interviewees (\texttt{CF1, CF2, CF8}), only 23\% of calculator users felt equipped to change \cite{collins2020}, and hostel, landlord-paid, and equal-split bills hide individual use (\texttt{CF1, CF7, CF10}). & \texttt{S1, L1, L3}; \texttt{P1, P3, P6}; \texttt{CF1} \cite{cf1}, \texttt{CF2} \cite{cf2}, \texttt{CF7} \cite{cf7}, \texttt{CF8} \cite{cf8}, \texttt{CF10} \cite{cf10}; LR \cite{lee2021,gandhi2025,grinstein2018,mulrow2019,collins2020} \\
\hline
\textbf{Pull}\newline (what is wanted) & Interviewees ask for terms they can picture (bus trips, 15L water buckets, fan/AC hourly costs, and rupees alongside kilos: \texttt{CF1, CF3, CF6, CF8}), visible Indian coal-grid assumptions (\texttt{CF4}), the easiest realistic swap (\texttt{CF5, CF6, CF10}), and a simple weekly or monthly view in regional languages (\texttt{CF1, CF3, CF4, CF7, CF8, CF10}) with a detailed view for experienced users (\texttt{CF4}, supported by \cite{hoffmann2024}). Users also respond to differences between options \cite{mohanty2023}. & \texttt{S1, S2, S3, L2}; \texttt{P1--P6}; \texttt{CF1} \cite{cf1}, \texttt{CF3}--\texttt{CF10} \cite{cf3,cf4,cf5,cf6,cf7,cf8,cf9,cf10}; LR \cite{hoffmann2024,mohanty2023} \\
\hline
\textbf{Anxiety}\newline (what puts users off) & Persuasive sustainability work has been criticised for narrowing sustainability to simple metrics, arguing against moralising scores \cite{brynjarsdottir2012}. Interviewees reject being lectured or made to feel like a bad person (\texttt{CF2, CF5, CF9}), distrust Western grid factors that underestimate India's coal grid (\texttt{CF4}), and drop 40-question forms, heavy typing, or data-hungry tools (\texttt{CF1, CF7}), matching evidence that tracking is effortful and leaves gaps when users stop logging \cite{hoffmann2024}. & \texttt{S1, S3, L3}; \texttt{P1--P4}; \texttt{CF1} \cite{cf1}, \texttt{CF2} \cite{cf2}, \texttt{CF4} \cite{cf4}, \texttt{CF5} \cite{cf5}, \texttt{CF7} \cite{cf7}, \texttt{CF9} \cite{cf9}; LR \cite{hoffmann2024,brynjarsdottir2012} \\
\hline
\textbf{Habit}\newline (existing routines and limits) & Footprints include emissions embodied in goods and services \cite{shi2021}, and tracking apps showed a weak mobility effect (12.0\% planned cut vs.\ 34.7\% for household activities) \cite{hoffmann2024}. Curfew cabs (\texttt{CF1}), night-shift cabs and daytime sleep AC (\texttt{CF9}), and 3 AM exam deliveries (\texttt{CF1}) are locked by safety, shift rotas, and heat. Low-waste routines (bucket baths, pressure cookers, repair, composting, \textit{kabadiwala} sales: \texttt{CF1, CF3, CF4, CF8}) go unrecognised, while one homemaker polices family showers and lights alone (\texttt{CF3}). & \texttt{S2, L1--L3}; \texttt{P1, P2, P4, P5, P6}; \texttt{CF1} \cite{cf1}, \texttt{CF3} \cite{cf3}, \texttt{CF4} \cite{cf4}, \texttt{CF5} \cite{cf5}, \texttt{CF8} \cite{cf8}, \texttt{CF9} \cite{cf9}; LR \cite{shi2021,hoffmann2024} \\
\hline
\end{longtable}
}

\subsubsection{Justification}
\textbf{``When I want to understand how my everyday routines and household habits affect the environment, I want to see which activities matter most...''}
The opening clause answers \textbf{Theme S2} and the scope gap in our literature review. Footprints vary across places and households in India \cite{lee2021} and include emissions embodied in goods and services, so estimates should state what they include \cite{shi2021}. The team’s early problem framing also highlighted three candidate areas for closer design attention—quick-commerce packaging and food waste, dairy-heavy food choices, and Tier-2-to-Tier-1 intercity travel (for example, bus-versus-rail comparisons). These are design opportunities identified for further validation, not yet proven to be the highest-emission activities for every user. The current interview themes directly support examining grocery delivery packaging, food waste, dairy and local alternatives; intercity travel should remain a candidate comparison scenario only if the reviewed literature or later requirements work can support a fair, sourced comparison. Most interviewed households eat vegetarian meals on most days or eat chicken only once or twice a week (\texttt{CF1, CF3, CF4, CF5, CF6, CF10}) \cite{cf1,cf3,cf4,cf5,cf6,cf10}. Their everyday concerns centre on food and grocery delivery plastic (\texttt{CF1, CF2, CF3, CF5, CF6, CF9}) \cite{cf1,cf2,cf3,cf5,cf6,cf9}, spoilt app vegetables and binned takeaway (\texttt{CF3, CF7}) \cite{cf3,cf7}, dairy versus local millets and rural-to-urban cooking fuels (\texttt{CF1, CF2, CF3, CF4, CF8}) \cite{cf1,cf2,cf3,cf4,cf8}, and one-season online clothes or annual phone upgrades (\texttt{CF2, CF10}) \cite{cf2,cf10}. At the same time, \texttt{P2}, \texttt{P5}, and \texttt{P6} (\texttt{SC1, SC4}) want existing low-waste habits recognised, including repairing phones, passing down clothes, bucket bathing, pressure-cooker cooking, composting, and selling dry waste to the \textit{kabadiwala} (\texttt{CF3, CF4, CF8}) \cite{cf3,cf4,cf8}.

\textbf{``...in terms I can picture... without being misled...''}
This part of the statement answers \textbf{Theme S1} and \textbf{Theme S3}, examined via \textbf{RQ1}. In a survey of 131 Delhi college students, only 38\% knew the term ``carbon footprint'' \cite{gandhi2025}, people estimate carbon quantities poorly relative to familiar metrics such as calories \cite{grinstein2018}, and calculators vary widely in regional customisation \cite{mulrow2019}. In our transcripts, \texttt{P1} (\texttt{CF1} \cite{cf1}), \texttt{P3} (\texttt{CF2} \cite{cf2}), \texttt{P5} (\texttt{CF6} \cite{cf6}), and \texttt{P6} (\texttt{CF8} \cite{cf8}) ask for everyday equivalents such as city bus trips, 15-litre water buckets, hourly fan or AC costs, and Indian rupees alongside weight, while \texttt{P2} (\texttt{CF4} \cite{cf4}) distrusts tools that hide Western grid assumptions instead of showing an Indian coal-grid factor. Because Mohanty et al.\ \cite{mohanty2023} found in hypothetical ride-hailing surveys that raw carbon weight outperformed relatable equivalencies and that users responded to differences between options rather than absolute values, which units work best remains unsettled. We therefore treat interviewees' requests for everyday household units alongside carbon weight as a stated preference to test through \texttt{RQ1} (\texttt{SC1, SC2, SC4}).

\textbf{``...and which realistic alternatives are open to me, so that I can make changes that fit my budget, safety, rota... without being... lectured, or blamed for constraints I cannot control.''}
Focusing on realistic alternatives and unavoidable limits answers \textbf{Latent Theme L3} and \textbf{Theme S3} (\texttt{RQ2}). Among 4,245 users of one widely used calculator, 91\% found it valuable for knowledge, yet only 23\% said it gave them the information needed to change their lives \cite{collins2020}. Hoffmann et al.\ \cite{hoffmann2024} found a weaker planned reduction for mobility (12.0\%) than for household activities (34.7\%), and Brynjarsd\'{o}ttir et al.\ \cite{brynjarsdottir2012} criticise persuasive sustainability work for narrowing sustainability to simple metrics, which argues against moralising scores. In our interviews, late-night cabs before a 10:30 PM hostel curfew (\texttt{P1} in \texttt{SC2}, \texttt{CF1} \cite{cf1}) or after a hospital night shift when buses have stopped (\texttt{P4} in \texttt{SC3}, \texttt{CF9} \cite{cf9}), daytime AC for post-shift sleep (\texttt{P4} in \texttt{SC3}, \texttt{CF9} \cite{cf9}), and 3 AM exam deliveries (\texttt{P1} in \texttt{SC2}, \texttt{CF1} \cite{cf1}) are locked by safety, shift rotas, and heat. Users also drop tools that demand 40-question forms, heavy typing, or mobile data drain (\texttt{CF1, CF7} \cite{cf1,cf7}; \cite{hoffmann2024}), asking instead for two-tap presets (\texttt{RQ2}), rota-aware quiet hours (\texttt{P4} in \texttt{SC3}), and one realistic swap or monthly challenge (\texttt{P3, P5} in \texttt{SC1, SC4}; \texttt{CF2, CF5, CF6, CF10} \cite{cf2,cf5,cf6,cf10}).

\textbf{``...and shared home without being... blamed...''}
Addressing the shared home without blame answers \textbf{Latent Themes L1 and L2} and \textbf{RQ3}. In \texttt{L1}, fixed hostel fees (\texttt{P1} in \texttt{SC2}, \texttt{CF1} \cite{cf1}), landlord-paid electricity (\texttt{CF7} \cite{cf7}), and equal bill splits in co-living flats (\texttt{P2} and \texttt{P3} in \texttt{SC4}, \texttt{CF10} \cite{cf10}) separate personal use from cost, so nobody wants to police the shared geyser or lounge AC. In \texttt{L2}, a single homemaker (\texttt{P5} in \texttt{SC1}, \texttt{CF3} \cite{cf3}) carries the tiring job of reminding teenagers about 20-minute showers and unattended lights while elders' (\texttt{P6} in \texttt{SC1}, \texttt{CF8} \cite{cf8}) bucket and fan habits are called ``old-fashioned''. Prior household water and shower eco-feedback studies \cite{froehlich2012,tiefenbeck2019} are formative and Western, with no Indian evidence on shared households, although real-time shower feedback cut hotel guests' energy use by 11.4\% even when they did not pay the bill \cite{tiefenbeck2019}. Across \texttt{P1} to \texttt{P6} and \texttt{SC1} to \texttt{SC4}, the statement serves all three stakeholder groups: urban students and shared-housing youth (\texttt{P1, P2}), time-poor working professionals (\texttt{P3, P4}), and household managers and elders (\texttt{P5, P6}).

% =========================================================================
% SECTION 3: PRIORITISATION METHOD
% =========================================================================
\subsection{Prioritisation Method}
Epics (\texttt{E1}--\texttt{E14}) are ranked with \textbf{MoSCoW} using explicit decision rules applied in order, aligned with our Java and JavaFX MVP scope (Section 5 of the Interview Analysis). A key scope distinction is intentional: displaying a transparent household total can be supported by the core estimate and assumptions journey, while richer coordination and voluntary per-profile attribution are kept as Should-Have until the team tests the consent model and confirms that inputs support meaningful attribution. This avoids presenting a shared bill as if it measured each person’s individual emissions. The rules are:
\begin{itemize}
    \item \textbf{Must have:} The core journey in at least two scenarios (\texttt{SC1}--\texttt{SC4}) cannot be completed without it, or omitting it would make the tool misleading, exclusionary, or judgmental toward safety- and shift-locked habits.
    \item \textbf{Should have:} It strengthens side-by-side comparison, repeat review, shared household action (\texttt{RQ3}), or researcher CSV export (Increment 2), while a usable first journey still works without it.
    \item \textbf{Could have:} It adds rota-aware quiet timing, static reference-flat benchmarking tables, or smart-plug timer helpers once Must and Should epics are in place.
    \item \textbf{Won't have (this version):} It is not required by \texttt{SC1}--\texttt{SC4}, or it needs automated bank/hardware data collection, forced roommate policing, or unverified carbon-offset accounting that cannot yet be justified. Each exclusion states a revisit condition.
\end{itemize}

% =========================================================================
% SECTION 4: EPICS AND MOSCOW PRIORITISATION
% =========================================================================
\subsection{Epics and MoSCoW Prioritisation}

{\footnotesize
\rowcolors{2}{bgRow}{white}
\begin{longtable}{|>{\raggedright\arraybackslash}p{2.7cm}|>{\raggedright\arraybackslash}p{5.2cm}|>{\raggedright\arraybackslash}p{2.9cm}|>{\raggedright\arraybackslash}p{5.2cm}|}
\hline
\rowcolor{cNavy}
\textcolor{white}{\textbf{Epic}} & \textcolor{white}{\textbf{High-Level User Story}} & \textcolor{white}{\textbf{Traceability}} & \textcolor{white}{\textbf{Acceptance Direction \& Rationale}} \\
\hline
\endfirsthead
\hline
\rowcolor{cNavy}
\textcolor{white}{\textbf{Epic}} & \textcolor{white}{\textbf{High-Level User Story}} & \textcolor{white}{\textbf{Traceability}} & \textcolor{white}{\textbf{Acceptance Direction \& Rationale}} \\
\hline
\endhead
\hline
\endfoot

\rowcolor{bgSand}\multicolumn{4}{|l|}{\textbf{\textcolor{cTerracotta}{MUST HAVE: The Core Journey and Its Trustworthiness}}} \\
\hline
\textbf{E1} Understand my footprint in familiar terms & As a user new to carbon footprints, I want to enter a manageable set of everyday activities and see the result in rupees, 15L water buckets, fan/AC hours, or bus trips beside \texttt{kg CO2e}, so that I can see which areas matter most to me. & \texttt{S1}; \texttt{RQ1}\newline \texttt{P1, P3, P5, P6}\newline \texttt{SC1, SC2, SC4} & A short 3-step setup finishes without requiring every field. Each result shows INR or a familiar equivalent beside \texttt{kg CO2e} \cite{mohanty2023,cf1,cf6,cf8}. Equivalents that show water volume or cost rather than carbon weight are clearly labelled. \\
\hline
\textbf{E2} See and edit the working behind an estimate & As a user who doubts unfamiliar or foreign numbers, I want to see and edit the assumed grid factor (\texttt{0.716 kg CO2/kWh}), tariff (\texttt{Rs.\ 8/kWh}), power draw, and limits behind each estimate, so that I can judge how far to trust it. & \texttt{S1, S3}\newline \texttt{P2, P3, P5, P6}\newline \texttt{SC1--SC4} & Every estimate opens a plain-language assumptions panel naming the factor and formula (\texttt{CF4} \cite{cf4}). Illustrative values are marked as assumed and editable in the MVP. \\
\hline
\textbf{E3} Find actions that fit my circumstances & As a user whose choices are limited by night safety curfews, rotating shifts, heat, or housing, I want suggestions that respect those limits, so that I get one realistic action rather than a lecture. & \texttt{L3, S3}\newline \texttt{P1, P3, P4}\newline \texttt{SC2, SC3, SC4} & Choices can be marked as locked necessities (curfew cab, night-shift cab, daytime sleep AC) and are never penalised \cite{brynjarsdottir2012,cf1,cf9}. Offers an achievable next step or single monthly challenge (\texttt{P3} \cite{cf2,cf10}). \\
\hline
\textbf{E4} See relevant Indian categories and circular habits & As a user whose impact extends beyond travel and electricity, I want delivery packaging, food waste, dairy vs.\ millets, cooking fuels, clothing/phone longevity, and repair or composting habits recorded, so that important activities are not hidden. & \texttt{S2}\newline \texttt{P1--P6}\newline \texttt{SC1--SC4} & Covered categories from \texttt{S2} are listed and excluded ones are stated \cite{shi2021,cf1,cf3,cf4,cf10}. Low-waste practices (bucket watering, pressure cooking, composting, \textit{kabadiwala} sales, cycling, phone repair) are recorded and visibly acknowledged (\texttt{SC1, SC4}). \\
\hline
\textbf{E5} Log briefly, offline, in my language & As a user with limited time, mobile data, or English, I want to record selected activities in 2 taps offline and read results in my language, so that I keep using the tool and can share it with family or flatmates. & \texttt{S3}; \texttt{RQ2}\newline \texttt{P1, P2, P4, P5, P6}\newline \texttt{SC1--SC4} & Core logging works 100\% offline with zero mobile data drain (\texttt{CF1} \cite{cf1}). Tap and time targets are tested via \texttt{RQ2} \cite{hoffmann2024,cf1,cf7}. English and Plain Hindi form the baseline (\texttt{SC1}), with Marathi, Tamil, Bengali, and Gujarati cards in \texttt{SC2} and \texttt{SC4}. \\
\hline

\rowcolor{bgSand}\multicolumn{4}{|l|}{\textbf{\textcolor{cTerracotta}{SHOULD HAVE: Comparison, Shared Household Action, Repeat Review, and Audit Export}}} \\
\hline
\textbf{E6} Compare feasible alternatives side by side & As a user choosing between options, I want to compare likely impact and cost (wasted geyser hours in bus trips, 1 hr less WFH AC, shower vs.\ buckets, or car vs.\ Metro and market vs.\ delivery), so that I can choose what suits me. & \texttt{S1, S2, L3}; \texttt{RQ1}\newline \texttt{P1, P3, P5}\newline \texttt{SC1, SC2, SC4} & Shows differences between options, since users respond to differences \cite{mohanty2023} (e.g., \texttt{2 hrs 2 kW geyser = 2.86 kg CO2 = 13 bus trips} in \texttt{SC2}; \texttt{1 hr less AC across 14 WFH days = 15.04 kg CO2 and Rs.\ 168/mo saved} in \texttt{SC4}; \texttt{20-min shower = ten 15L buckets} in \texttt{SC1}). \\
\hline
\textbf{E7} Coordinate shared household actions & As a person sharing a family home, hostel floor, or co-living flat, I want shared loads shown as household totals by default while letting each member voluntarily opt in on their own profile to showing their individual share, so that nobody has to police the rest. & \texttt{L1, L2}; \texttt{RQ3}\newline \texttt{P1, P2, P3, P5, P6}\newline \texttt{SC1, SC2, SC4} & Shows shared household totals before individual shares (\texttt{RQ3}) \cite{froehlich2012,tiefenbeck2019}. Per-profile opt-in (\texttt{SC1, SC4}) is strictly voluntary: no member can switch attribution on for others, and non-opted-in profiles show no pressure prompt. \\
\hline
\textbf{E8} Review change over time and monthly challenges & As a user who has picked a habit swap or monthly challenge (such as keeping a phone 1 extra year), I want a simple weekly or monthly summary of earlier records, so that I can see whether my habits have changed. & \texttt{S2, S3}\newline \texttt{P1, P3, P5}\newline \texttt{SC2, SC4} & Past records can be viewed in a simple weekly/monthly picture (\texttt{CF2, CF7, CF10} \cite{cf2,cf7,cf10}) with no streaks or scores, as reviewed studies do not show that streaks sustain use \cite{hoffmann2024,froehlich2010}. \\
\hline
\textbf{E9} Open a Detailed Audit View and export CSV data & As a user who wants to verify the evidence, I want a Detailed Audit View of factors and formulas and CSV data export (Increment 2) while the Simple View stays default, so that I can audit results for research. & \texttt{S1, S3}\newline \texttt{P2} (Simple View for \texttt{P1, P6})\newline \texttt{SC4} (Inc.\ 2) & Simple View stays default for newcomers \cite{hoffmann2024}. Detailed Audit View and CSV export are scoped as Should-Have (Increment 2) to match Section 5, \texttt{P2}, and \texttt{SC4} of TA v1.7 \cite{cf4}. \\
\hline

\rowcolor{bgSand}\multicolumn{4}{|l|}{\textbf{\textcolor{cTerracotta}{COULD HAVE: Rota Timing, Static Reference-Flat Benchmarks, and Timer Helpers}}} \\
\hline
\textbf{E10} Control rota-aware feedback timing & As a rotating night-shift worker, I want to enter my weekly rota in 2 taps so daytime sleep alerts are muted and I get one short weekly message on a quiet rest day, so that feedback never disturbs my sleep. & \texttt{L3, S3}\newline \texttt{P4}\newline \texttt{SC3} & No unsolicited notifications by default. Entering a weekly rota (\texttt{P4} in \texttt{SC3} \cite{cf9}) mutes daytime sleep hours and schedules one weekly rest-day prompt (washing load and steel tiffin prep). \\
\hline
\textbf{E11} Benchmark against a reference flat table & As a user surprised by an electricity bill, I want to see my usage beside a peer or reference-flat benchmark using a static table, so that I know whether my load is typical. & \texttt{S1}\newline \texttt{P3}\newline \texttt{SC4} & Scoped as a Could-Have using a static reference-flat lookup table (matching Section 5, \texttt{P3}, and \texttt{SC4} of TA v1.7 \cite{lee2021,cf2}), stating all reference flat assumptions clearly. \\
\hline
\textbf{E12} Record smart-plug cut-off timers & As a user who already owns a basic smart plug for a geyser or room AC, I want to record a scheduled cut-off timer (such as turning the AC off 1 hour earlier), so that I can track savings without daily effort. & \texttt{S3, L2}\newline \texttt{P3, P5}\newline \texttt{SC1, SC4} & Records a user-set timer schedule for an existing smart plug (\texttt{P5} morning geyser timer in \texttt{SC1} \cite{cf3}; \texttt{P3} WFH AC cut-off in \texttt{SC4} \cite{cf2}) without needing external cloud hardware APIs. \\
\hline

\rowcolor{bgSand}\multicolumn{4}{|l|}{\textbf{\textcolor{cTerracotta}{WON'T HAVE (This Version): Excluded with a Revisit Condition}}} \\
\hline
\textbf{E13} Public rankings and forced per-person bill splits & Users will not be ranked on leaderboards or have non-consensual per-person shares of a shared bill switched on for them by another household member. & \texttt{L1, L2, L3}; \texttt{RQ3}\newline \texttt{P1, P2, P3, P5}\newline Excluded in \texttt{SC1, SC2, SC4} & Rankings ignore unequal safety and shift constraints \cite{brynjarsdottir2012}, and shared bills do not measure individual fixture use \cite{froehlich2012}. Revisit only with fixture-level sub-metering and explicit consent. \\
\hline
\textbf{E14} Automatic hardware/bank scraping and unverified offset credits & Users will not have automatic inverter battery telemetry (\texttt{CF3, CF8}), bank/order statement scraping (\texttt{CF2}), building construction LCA (\texttt{CF4}), or unverified carbon-offset numbers in v1.0. & \texttt{S2, S3}\newline \texttt{P2, P3, P5, P6}\newline Scoped out for v1.0 & Follows \texttt{CF4}'s warning (\textit{``that's probably too much for a first version. Don't overbuild it''} \cite{cf4}). Circular habits are credited descriptively on the dashboard (\texttt{E4}) until a sourced offset method is chosen. \\
\hline
\end{longtable}
}

% =========================================================================
% SECTION 5: SCENARIO AND PERSONA TRACEABILITY
% =========================================================================
\subsection{Scenario and Persona Traceability}

\subsubsection{Scenario-to-Epic Mapping (Synchronised with TA v1.7 Section 7)}
{\footnotesize
\rowcolors{2}{bgRow}{white}
\begin{longtable}{|>{\raggedright\arraybackslash}p{5.3cm}|>{\raggedright\arraybackslash}p{2.8cm}|>{\raggedright\arraybackslash}p{2.3cm}|>{\raggedright\arraybackslash}p{5.6cm}|}
\hline
\rowcolor{cNavy}
\textcolor{white}{\textbf{Scenario (TA v1.7 Section 7)}} & \textcolor{white}{\textbf{Personas}} & \textcolor{white}{\textbf{Must Epics}} & \textcolor{white}{\textbf{Supporting Epics (\& Exclusions)}} \\
\hline
\textbf{SC1:} Resolving Summer Tanker-Water and Fan/Geyser Friction in a Noida 3BHK Home & \texttt{P5} Sunita, \texttt{P6} Rajeshwar, \& Teenagers & \texttt{E1, E2, E4, E5} & \texttt{E6} (20-min shower = ten 15L buckets; fan Rs.\ 0.60/hr), \texttt{E7} (shared total default + voluntary per-profile opt-in), \texttt{E12} (geyser timer); \texttt{E13, E14} excluded \\
\hline
\textbf{SC2:} Late-Night Exam Delivery Logging and Shared Geyser Nudges in a Pune Hostel & \texttt{P1} Ananya \& Hostel Floor-Mates & \texttt{E1, E2, E3, E4, E5} & \texttt{E6} (2 hrs 2 kW geyser = 2.86 kg CO$_2$ = 13 bus trips), \texttt{E7} (Marathi/Hindi/Bengali shared floor card), \texttt{E8} (weekly picture); \texttt{E13} excluded \\
\hline
\textbf{SC3:} Night-Shift Rota Quiet Hours and Rest-Day Planning in a Hyderabad Flat & \texttt{P4} Priya & \texttt{E2, E3, E4, E5} & \texttt{E10} (2-tap rota entry, daytime sleep alert muting, one weekly rest-day washing \& steel tiffin prompt) \\
\hline
\textbf{SC4:} Shared Co-Living Bill Audit, WFH AC, and Phone Longevity in a Chennai 2BHK & \texttt{P2} Karthik \& \texttt{P3} Rohan & \texttt{E1, E2, E3, E4, E5} & \texttt{E6} (1 hr earlier AC cut-off = 15.04 kg CO$_2$ \& Rs.\ 168/mo saved), \texttt{E7} (voluntary room opt-in), \texttt{E8} (keep phone 1 extra year), \texttt{E9} (Detailed Audit View + CSV export, Inc.\ 2), \texttt{E11} (static reference-flat table), \texttt{E12} (AC timer); \texttt{E13, E14} excluded \\
\hline
\end{longtable}
}

\subsubsection{Persona Coverage Check}
Every persona (\texttt{P1}--\texttt{P6}) is served by at least four Must Have epics and supported by Should or Could epics: \textbf{P1 Ananya} (\texttt{E1--E8} in \texttt{SC2}), \textbf{P2 Karthik} (\texttt{E1, E2, E4, E5, E7, E9} in \texttt{SC4}), \textbf{P3 Rohan} (\texttt{E1--E4, E6--E8, E11, E12} in \texttt{SC4}), \textbf{P4 Priya} (\texttt{E2--E5, E10} in \texttt{SC3}), \textbf{P5 Sunita} (\texttt{E1, E2, E4--E8, E12} in \texttt{SC1}), and \textbf{P6 Rajeshwar} (\texttt{E1, E2, E4--E7} in \texttt{SC1}).

% =========================================================================
% SECTION 6: WHY THESE PRIORITIES? TENSIONS AND CONSISTENCY CHECKS
% =========================================================================
\subsection{Why These Priorities? Tensions, Decisions, and Consistency Checks}

The Must epics (\texttt{E1}--\texttt{E5}) form the minimum trustworthy loop: understand impact in familiar terms (\texttt{E1}), inspect and edit calculation assumptions (\texttt{E2}), respect safety, shift-work, and heat constraints without moralising scores (\texttt{E3}), see relevant Indian convenience and circular habits (\texttt{E4}), and log offline in two taps across local languages (\texttt{E5}). Five tensions from our triangulation table shaped the ranking:
\begin{enumerate}
    \item \textbf{T1: Familiar Equivalents vs.\ Carbon Weight (\texttt{E1, E6} | \texttt{RQ1}):} Interviewees ask for water buckets, fan-hours, bus trips, and rupees \cite{cf1,cf6,cf8}, whereas raw $\text{CO}_2$ weight outperformed equivalencies in hypothetical ride-hailing choices and users responded to differences between options \cite{mohanty2023}. \textit{Decision:} Show everyday units and INR beside \texttt{kg CO2e} (\texttt{E1}), highlight differences between options (\texttt{E6}), and test comprehension via \texttt{RQ1}.
    \item \textbf{T2: Seeing ``Their Own Part'' vs.\ Avoiding Unsupported Attribution (\texttt{E7, E13} | \texttt{RQ3}):} \texttt{CF3} asks for family members to \textit{``see their own part''} \cite{cf3}, while LR Table 1 specifies \textit{``Shared totals before individual shares; no unsupported attribution''} and \texttt{RQ3} asks how a shared view can show communal totals without allocating them to individuals. \textit{Decision (matching TA v1.7 Section 5, SC1, and SC4):} Show shared household totals by default (\texttt{E7}) and let each member voluntarily opt in on their own profile to showing their individual or room share where inputs support it, while forbidding anyone from switching attribution on for others (\texttt{E13}).
    \item \textbf{T3: Recognition of Repair and Composting vs.\ Accounting Rigour (\texttt{E4, E14}):} \texttt{P2, P5}, and \texttt{P6} want bucket watering, pressure cooking, phone repair, composting, and \textit{kabadiwala} sales recognised \cite{cf3,cf4,cf8}. \textit{Decision:} Record and visibly credit these practices on the dashboard (\texttt{E4}) while deferring unverified carbon-offset numbers (\texttt{E14}).
    \item \textbf{T4: Audit Depth for Experts vs.\ Simplicity for Newcomers and Elders (\texttt{E2, E5, E9} | \texttt{RQ2}):} \texttt{P2} wants to check grid factors and export CSV records \cite{cf4}, while \texttt{P1} and \texttt{P6} need simple, jargon-free screens \cite{cf1,cf8}. \textit{Decision:} Keep the 2-tap Simple View as the default (\texttt{E5}) with an always-available assumptions panel (\texttt{E2}), and deliver the Detailed Audit View and CSV export in Increment 2 (\texttt{E9}, Should-Have) \cite{hoffmann2024}.
    \item \textbf{T5: Protecting Shift-Worker Rest vs.\ Keeping the Core Small (\texttt{E3, E10, E12}):} \texttt{P4} needs quiet daytime sleep windows around her hospital rota \cite{cf9}. \textit{Decision:} Send no unsolicited notifications by default (\texttt{E10}, Could-Have) and never penalise safety- or health-locked choices in the core (\texttt{E3}).
\end{enumerate}

\begin{checkcard}{Explicit Consistency \& Alignment Notes (LR v1.0 and TA v1.7)}
\footnotesize
\begin{itemize}
    \item \textbf{Personas \& Scenarios Verified (TA v1.7 Sections 6 \& 7):} Checked directly against TA v1.7: \textbf{SC1} (\texttt{P5} Sunita, \texttt{P6} Rajeshwar relocated to Noida, and Sunita's two teenagers), \textbf{SC2} (\texttt{P1} Ananya and Pune hostel floor-mates), \textbf{SC3} (\texttt{P4} Priya in Hyderabad), and \textbf{SC4} (\texttt{P2} Karthik and composite persona \texttt{P3} Rohan relocated to share a Chennai 2BHK).
    \item \textbf{Flagged Differences Between LR v1.0 and TA v1.7 (Not Silently Resolved):} (1) \textit{Attribution (\texttt{RQ3}):} LR Section 5 words \texttt{RQ3} as \textit{``without allocating them to individuals''}, while TA v1.7 Section 5 (p.\ 7) proposes shared totals by default with voluntary per-profile opt-in (\texttt{SC1, SC4}), noting that this goes slightly beyond \texttt{RQ3} as worded and will be settled in D2; (2) \textit{Stakeholder Scope:} LR focuses on urban students and young adults (\texttt{P1--P4}), while TA v1.7 adds family personas (\texttt{P5, P6}) to cover transcripts \texttt{CF3, CF6, CF8}; (3) \textit{Units \& Placeholders:} LR uses \texttt{kg CO2e} while TA v1.7 scenario calculations use \texttt{kg CO2} with illustrative, unsourced assumptions (\texttt{0.716 kg CO2/kWh}, \texttt{0.015 kg CO2/pkm}, \texttt{Rs.\ 8/kWh}) marked with \texttt{[VERIFY]} placeholders in TA v1.7 Section 7.
\end{itemize}
\end{checkcard}

% =========================================================================
% SECTION 7: LIMITATIONS AND HAND-OFF TO D2-D6
% =========================================================================
\subsection{Limitations and Hand-Off to D2--D6}
Our priorities rest on ten illustrative coursework transcripts (\texttt{CF1}--\texttt{CF10}) \cite{cf1,cf10} and on a literature review that identifies gaps in the sources reviewed rather than the whole literature. They are therefore testable hypotheses about stakeholder value. In Deliverables 2 to 6, we will (1) split \texttt{E1}--\texttt{E12} into granular user stories with measurable acceptance criteria, (2) settle the MVP category subset, language set, and tap/time targets, (3) evaluate \texttt{RQ1} (units, \texttt{E1/E6}), \texttt{RQ2} (2-tap presets, \texttt{E5}), and \texttt{RQ3} (shared household views, \texttt{E7}), and (4) revisit \texttt{E13} and \texttt{E14} only when their stated conditions are met.

% =========================================================================
% REFERENCES (EXACT NUMBERING [1]-[23] FROM LR v1.0 & TA v1.7)
% =========================================================================

\clearpage
\begin{thebibliography}{99}
\bibitem{lee2021}
J. Lee, O. Taherzadeh, and K. Kanemoto. 2021. The scale and drivers of carbon footprints in households, cities and regions across India. \emph{Global Environmental Change} 66, 102205. \doi{10.1016/j.gloenvcha.2020.102205}

\bibitem{gandhi2025}
N. Gandhi and P. Shukla. 2025. Dietary sustainability knowledge and practice among college-going students in Delhi. \emph{Journal of Nutrition and Food Security} 10(4), 583--593. \doi{10.18502/jnfs.v10i4.20014}

\bibitem{grinstein2018}
A. Grinstein, E. Kodra, S. Chen, S. Sheldon, and O. Zik. 2018. Carbon innumeracy. \emph{PLOS ONE} 13(5), e0196282. \doi{10.1371/journal.pone.0196282}

\bibitem{shi2021}
S. Shi and J. Yin. 2021. Global research on carbon footprint: A scientometric review. \emph{Environmental Impact Assessment Review} 89, 106571. \doi{10.1016/j.eiar.2021.106571}

\bibitem{mulrow2019}
J. Mulrow, K. Machaj, J. Deanes, and S. Derrible. 2019. The state of carbon footprint calculators: An evaluation of calculator design and user interaction features. \emph{Sustainable Production and Consumption} 18, 33--40. \doi{10.1016/j.spc.2018.12.001}

\bibitem{collins2020}
A. Collins, A. Galli, T. Hipwood, and A. Murthy. 2020. Living within a One Planet reality: The contribution of personal Footprint calculators. \emph{Environmental Research Letters} 15, 025008. \doi{10.1088/1748-9326/ab5f96}

\bibitem{hoffmann2024}
S. Hoffmann, W. Lasarov, H. Reimers, and M. Trabandt. 2024. Carbon footprint tracking apps. Does feedback help reduce carbon emissions? \emph{Journal of Cleaner Production} 434, 139981. \doi{10.1016/j.jclepro.2023.139981}

\bibitem{lasarov2024}
W. Lasarov, S. Hoffmann, R. Mai, and J. Schleich. 2024. Carbon footprint tracking apps: The spillover effects of feedback and goal-activating appeals. \emph{Ecological Economics} 226, 108354. \doi{10.1016/j.ecolecon.2024.108354}

\bibitem{froehlich2010}
J. Froehlich, L. Findlater, and J. Landay. 2010. The design of eco-feedback technology. In \emph{Proceedings of the Conference on Human Factors in Computing Systems (CHI '10)}, 1999--2008. Association for Computing Machinery (ACM). \doi{10.1145/1753326.1753629}

\bibitem{mohanty2023}
V. Mohanty, A. Filipowicz, N. Bravo, S. Carter, and D. A. Shamma. 2023. Save a tree or 6 kg of CO$_2$? Understanding effective carbon footprint interventions for eco-friendly vehicular choices. In \emph{Proceedings of CHI '23}. ACM. \doi{10.1145/3544548.3580675}

\bibitem{brynjarsdottir2012}
H. Brynjarsd\'ottir, M. H\r{a}kansson, J. Pierce, E. Baumer, C. DiSalvo, and P. Sengers. 2012. Sustainably unpersuaded: How persuasion narrows our vision of sustainability. In \emph{Proceedings of CHI '12}, 947--956. ACM. \doi{10.1145/2207676.2208539}

\bibitem{froehlich2012}
J. Froehlich, L. Findlater, M. Ostergren, et al. 2012. The design and evaluation of prototype eco-feedback displays for fixture-level water usage data. In \emph{Proceedings of CHI '12}, 2367--2376. ACM. \doi{10.1145/2207676.2208397}

\bibitem{tiefenbeck2019}
V. Tiefenbeck, A. W\"orner, S. Sch\"ob, E. Fleisch, and T. Staake. 2019. Real-time feedback promotes energy conservation in the absence of volunteer selection bias and monetary incentives. \emph{Nature Energy} 4(1), 35--41. \doi{10.1038/s41560-018-0282-1}

\bibitem{cf1}
COMP2300 Module Team. 2026. \textit{Interview Transcript CF1: Woman, Age 21, Final-Year Engineering Student, Hostel in Pune}. In COMP2300 Transcripts: Carbon Footprint in India (realistic v2), University of Southampton, 2026-10-19.

\bibitem{cf2}
COMP2300 Module Team. 2026. \textit{Interview Transcript CF2: Man, Age 29, Hybrid Software Engineer, Lives Alone in Bengaluru}. In COMP2300 Transcripts: Carbon Footprint in India (realistic v2), University of Southampton, 2026-10-19.

\bibitem{cf3}
COMP2300 Module Team. 2026. \textit{Interview Transcript CF3: Woman, Age 46, Homemaker and Part-Time Accounts Assistant, 3BHK in Noida}. In COMP2300 Transcripts: Carbon Footprint in India (realistic v2), University of Southampton, 2026-10-19.

\bibitem{cf4}
COMP2300 Module Team. 2026. \textit{Interview Transcript CF4: Man, Age 24, Environmental Science Postgraduate, PG in Chennai (Village Roots in Tamil Nadu)}. In COMP2300 Transcripts: Carbon Footprint in India (realistic v2), University of Southampton, 2026-10-19.

\bibitem{cf5}
COMP2300 Module Team. 2026. \textit{Interview Transcript CF5: Woman, Age 34, Marketing Manager, Lives with Partner in Thane, Commutes to Andheri}. In COMP2300 Transcripts: Carbon Footprint in India (realistic v2), University of Southampton, 2026-10-19.

\bibitem{cf6}
COMP2300 Module Team. 2026. \textit{Interview Transcript CF6: Man, Age 38, Bank Officer, Lives with Wife and Daughter in Gurugram}. In COMP2300 Transcripts: Carbon Footprint in India (realistic v2), University of Southampton, 2026-10-19.

\bibitem{cf7}
COMP2300 Module Team. 2026. \textit{Interview Transcript CF7: Woman, Age 19, First-Year Undergraduate, Shares Flat with Two Roommates in Kolkata}. In COMP2300 Transcripts: Carbon Footprint in India (realistic v2), University of Southampton, 2026-10-19.

\bibitem{cf8}
COMP2300 Module Team. 2026. \textit{Interview Transcript CF8: Man, Age 58, Retired Railway Official, Lives with Wife and Grandson in Lucknow}. In COMP2300 Transcripts: Carbon Footprint in India (realistic v2), University of Southampton, 2026-10-19.

\bibitem{cf9}
COMP2300 Module Team. 2026. \textit{Interview Transcript CF9: Woman, Age 31, Staff Nurse on Rotating Night Shifts, Lives Alone in Hyderabad}. In COMP2300 Transcripts: Carbon Footprint in India (realistic v2), University of Southampton, 2026-10-19.

\bibitem{cf10}
COMP2300 Module Team. 2026. \textit{Interview Transcript CF10: Man, Age 26, Freelance Graphic Designer, Co-Living Space in Ahmedabad}. In COMP2300 Transcripts: Carbon Footprint in India (realistic v2), University of Southampton, 2026-10-19.
\end{thebibliography}
\end{document}